\documentclass[JEP,XML,SOM,Unicode,NoEqCountersInSection,published]{cedram}
\datereceived{2025-02-14}
\dateaccepted{2025-07-30}
\dateepreuves{2025-08-25}

\usepackage[scr=rsfs,cal=euler]{mathalfa}
\renewcommand{\thetable}{\arabic{table}}
\multlinegap0pt

\hyphenation{every-thing recall table}

\usepackage[capitalise]{cleveref}
\crefdefaultlabelformat{#2\textup{#1}#3}

\makeatletter
\def\cref@thmoptarg[#1]#2#3#4{%
    \ifhmode\unskip\unskip\par\fi%
    \normalfont%
    \trivlist%
    \let\thmheadnl\relax%
    \let\thm@swap\@gobble%
    \thm@notefont{\fontseries\mddefault\upshape}%
    \thm@headpunct{.}% add period after heading
    \thm@headsep 5\p@ plus\p@ minus\p@\relax%
    \thm@space@setup%
    #2% style overrides
    \@topsep \thm@preskip               % used by thm head
    \@topsepadd \thm@postskip           % used by \@endparenv
    \def\@tempa{#3}\ifx\@empty\@tempa%
      \def\@tempa{\@oparg{\@begintheorem{#4}{}}[]}%
    \else%
      \refstepcounter[#1]{#3}%  <<< cleveref modification
      \@namedef{cref@#3@alias}{#1}% added
      \def\@tempa{\@oparg{\@begintheorem{#4}{\csname the#3\endcsname}}[]}%
    \fi%
    \@tempa}%
\makeatother

\let\oldbigsqcup\bigsqcup
\renewcommand{\bigsqcup}{\mathop{\textstyle\oldbigsqcup}\displaylimits}
\newcommand\mto{\mathchoice{\longmapsto}{\mapsto}{\mapsto}{\mapsto}}
\newcommand{\sbullet}{{\scriptscriptstyle\bullet}}
\newcommand{\csbullet}{{\raisebox{1pt}{$\sbullet$}}}

\let\epsilon\varepsilon
\let\ra\rightarrow

\usepackage{booktabs}
\usepackage{colonequals}
\usepackage{tikz-cd}

\theoremstyle{plain}
\newtheorem{theorem}{Theorem}[section]
\newtheorem{corollary}[theorem]{Corollary}
\newtheorem{lemma}[theorem]{Lemma}
\newtheorem{proposition}[theorem]{Proposition}

\newtheorem{alphatheorem}{Theorem}
\newtheorem{alphaconjecture}[alphatheorem]{Conjecture}
\newtheorem{alphacorollary}[alphatheorem]{Corollary}
\newtheorem{alphaproposition}[alphatheorem]{Proposition}

\theoremstyle{definition}
\newtheorem{definition}[theorem]{Definition}
\newtheorem{remark}[theorem]{Remark}
\newtheorem{example}[theorem]{Example}

\renewcommand{\thealphatheorem}{\Alph{alphatheorem}}
\renewcommand{\thealphaconjecture}{\Alph{alphatheorem}}
\renewcommand{\thealphacorollary}{\Alph{alphatheorem}}
\renewcommand{\thealphaproposition}{\Alph{alphaproposition}}

\crefname{alphatheorem}{Theorem}{Theorems}
\crefname{alphaconjecture}{Conjecture}{Conjectures}
\crefname{alphacorollary}{Corollary}{Corollaries}
\crefname{alphaproposition}{Proposition}{Propositions}

\DeclareMathOperator\Br{Br}
\DeclareMathOperator\characters{X}
\DeclareMathOperator\derived{\mathbf{D}}
\DeclareMathOperator\dimvect{\underline{\operatorname{dim}}}
\DeclareMathOperator\End{End}
\DeclareMathOperator\Ext{Ext}
\DeclareMathOperator\GL{GL}
\DeclareMathOperator\HH{H}
\DeclareMathOperator\hochschild{HH}
\DeclareMathOperator\Hilb{Hilb}
\DeclareMathOperator\Hom{Hom}
\DeclareMathOperator\sheafHom{\mathcal{H}om}
\DeclareMathOperator\RsheafHom{\mathbf{R}\mathcal{H}om}
\DeclareMathOperator\Mat{Mat}
\DeclareMathOperator\PG{PG}
\DeclareMathOperator\Pic{Pic}
\DeclareMathOperator\Rep{R}
\DeclareMathOperator\rk{rk}
\DeclareMathOperator\Spec{Spec}
\DeclareMathOperator\source{s}
\DeclareMathOperator\target{t}
\DeclareMathOperator\weight{wt}

\mathchardef\mhyphen="2D

\newcommand\blade[1][]{\Sigma_{#1}}
\newcommand\bounded{\ensuremath{\mathrm{b}}}
\newcommand\can{\ensuremath{\mathrm{can}}}
\newcommand\cocharacters{\ensuremath{\mathrm{X}_*}}
\newcommand\dash{\nobreakdash-\hspace{0pt}}
\newcommand\diff{\mathrm{d} \kern 0pt}
\newcommand\field{\mathbf{k}}
\newcommand\Gm{\ensuremath{\mathbf{G}_{\mathrm{m}}}}
\newcommand\group[1][\tuple{d}]{\mathrm{G}_{#1}}
\newcommand\HN{\ensuremath{{\scriptscriptstyle\mathrm{HN}}}}
\newcommand\limitset[1][\lambda]{Z_{#1}}
\newcommand\moduli{\ensuremath{\mathrm{M}}}
\newcommand\normal{\ensuremath{\mathrm{N}}}
\newcommand\semisimple{\ensuremath{\mathrm{ssimp}}}
\newcommand\semistable{\ensuremath{\mhyphen\mathrm{sst}}}
\newcommand\stable{\ensuremath{\mhyphen\mathrm{st}}}
\newcommand\tangent{\ensuremath{\mathrm{T}}}
\newcommand\tuple[1]{\ensuremath{{\boldsymbol{#1}}}}

\DeclareDocumentCommand\modulistack{om}{\IfNoValueTF{#1}{\mathcal{M}{(#2)}}{\mathcal{M}^{#1}(#2)}}
\DeclareDocumentCommand\modulispace{om}{\IfNoValueTF{#1}{\mathrm{M}{(#2)}}{\mathrm{M}^{#1}(#2)}}
\DeclareDocumentCommand\repspace{om}{\IfNoValueTF{#1}{\mathrm{R}{(#2)}}{\mathrm{R}^{#1}(#2)}}

\newif\ifcode

\datepublished{2025-09-04}
\begin{document}
\frontmatter

\title[Rigidity and Schofield's partial tilting conjecture for quiver moduli]{\hbox{Rigidity} and Schofield's partial tilting conjecture for quiver moduli}

\author[\initial{P.} \lastname{Belmans}]{\firstname{Pieter} \lastname{Belmans}}
\address{Mathematical Institute, Utrecht University,\\
Budapestlaan 6, 3584 CD Utrecht, The Netherlands}
\email{p.belmans@uu.nl}
%\urladdr{https://pbelmans.ncag.info/}

\author[\initial{A.-M.} \lastname{Brecan}]{\firstname{Ana-Maria} \lastname{Brecan}}
\email{anabrecan@gmail.com}
\urladdr{}

\author[\initial{H.} \lastname{Franzen}]{\firstname{Hans} \lastname{Franzen}}
\address{Institute of Mathematics, Paderborn University,\\
Warburger Str.\ 100, 33098 Paderborn, Germany}
\email{hans.franzen@math.upb.de}
%\urladdr{https://math.uni-paderborn.de/ag/arbeitsgruppe-algebra/team/hans-franzen}

\author[\initial{G.} \lastname{Petrella}]{\firstname{Gianni} \lastname{Petrella}}
\address{Department of Mathematics, Université de Luxembourg,\\
2 Av.\ de l'Université, 4365, Esch-sur-Alzette, Luxembourg}
\email{gianni.petrella@uni.lu}
%\urladdr{https://giannipetrella.eu/}

\author[\initial{M.} \lastname{Reineke}]{\firstname{Markus} \lastname{Reineke}}
\address{Fakultat für Mathematik, Ruhr-Universität Bochum,\\ Universitätsstr.\ 150, 44780 Bochum, Germany}
\email{markus.reineke@rub.de}
%\urladdr{https://math.ruhr-uni-bochum.de/fakultaet/arbeitsbereiche/algebra/research-team-reineke/team/prof-dr-markus-reineke/}

\thanks{P.B. was partially supported by the Luxembourg National Research Fund (FNR--17113194).
H.F. was partially supported by the Deutsche Forschungsgemeinschaft (DFG, German Research Foundation) SFB-TRR~358/1~2023 ``Integral Structures in Geometry and Representation Theory'' (491392403).
G.P. was supported by the Luxembourg National Research Fund (FNR--17953441).
M.R. was supported by the Deutsche Forschungsgemeinschaft (DFG, German Research Foundation) CRC-TRR~191 ``Symplectic structures in geometry, algebra and dynamics'' (281071066)}

\begin{abstract}
We explain how Teleman quantization can be applied to moduli spaces of quiver representations, in order to compute the higher cohomology of the endomorphism bundle of the universal bundle. We use this to prove Schofield's partial tilting conjecture in many interesting cases, and to show that moduli spaces of quiver representations are often (infinitesimally) rigid as varieties.
\end{abstract}

\subjclass{14D20, 14F08, 16G20}

\keywords{Moduli spaces of quiver representations, Teleman quantization, deformation theory}

\altkeywords{Espaces de modules de représentations de carquois, quantification de Teleman, rigidité}

\alttitle{Rigidité et conjecture de tilting partiel de Schofield pour les espaces de modules de carquois}

\begin{altabstract}
Nous expliquons comment la quantification de Teleman peut être appliquée aux espaces de modules de représentations de carquois, afin de calculer la cohomologie supérieure du fibré des endomorphismes du fibré universel. Nous utilisons cela pour démontrer la conjecture de tilting partiel de Schofield dans de nombreux cas intéressants, et pour montrer que les espaces de modules de représentations de carquois sont souvent (infinitésimalement) rigides en tant que variétés.
\end{altabstract}

\maketitle

\tableofcontents
\mainmatter

\section{Introduction}
\label{section:introduction}
There exists a rich moduli theory for representations of a quiver~$Q=(Q_0,Q_1)$,
surveyed in~\cite{MR2484736}.
Following King~\cite{MR1315461},
one fixes a dimension vector~$\tuple{d}\in\mathbb{N}^{Q_0}$
and a stability parameter~$\theta\in\Hom(\mathbb{Z}^{Q_0},\mathbb{Z})$,
and constructs a moduli space~$\moduli^{\theta\stable}(Q,\tuple{d})$ as a GIT quotient,
parametrizing isomorphism classes of~$\theta$-stable representations
of dimension vector~$\tuple{d}$.

When~$\tuple{d}$ and~$\theta$ are chosen appropriately,
the moduli space~$\moduli^{\theta\stable}(Q,\tuple{d})$
is a \emph{fine} moduli space,
parametrizing isomorphism classes of~$\theta$-(semi)stable representations. It~comes with a universal representation~$\mathcal{U}$:
a vector bundle equipped with a~$\field Q$\nobreakdash-action,
whose fiber in a point~$[V]\in\moduli^{\theta\stable}(Q,\tuple{d})$
is the representation~$V$.
The universal bundle~$\mathcal{U}$ decomposes as a direct sum~$\bigoplus_{i\in Q_0}\mathcal{U}_i$.
These constructions and results will be recalled in \cref{section:moduli}.

In this paper
we are interested in~$\sheafHom(\mathcal{U},\mathcal{U})\cong\mathcal{U}^\vee\otimes\mathcal{U}\cong\bigoplus_{i,j\in Q_0}\mathcal{U}_i^\vee\otimes\mathcal{U}_j$,
and we prove the following result on the summands,
which can be equivalently stated by saying that~$\sheafHom(\mathcal{U},\mathcal{U})$
has no higher cohomology.
The condition of being strongly amply~$\theta$-stable
is given in \cref{definition:strong-ample-stability}.

\begin{alphatheorem}
\label{theorem:cohomology-vanishing}
Let~$Q$ be a quiver,
let~$\tuple{d}$ be a dimension vector,
and let~$\theta$ be a stability parameter
such that~$\tuple{d}$ is~$\theta$\dash coprime.
Assume furthermore that~$\tuple{d}$ is strongly amply~$\theta$\dash stable.\footnote{See \cref{remark:strong-ample-stability} for a weaker condition,
and \cref{remark:ample-stability} for speculations on a necessary condition.}
Then, for all~$i,j \in Q_0$ we have
\begin{equation}
\label{equation:cohomology-vanishing}
\HH^{\geq1}(\moduli^{\theta\stable}(Q,\tuple{d}), \mathcal{U}^\vee_i \otimes \mathcal{U}_j) = 0,
\end{equation}
where the $\mathcal{U}_i$ are the summands of
the universal representation on the moduli space.
\end{alphatheorem}

For this result we do not have to assume that~$Q$ is acyclic.
We will prove this using \emph{Teleman quantization},
a tool that allows to compute sheaf cohomology on a GIT~quotient
using sheaf cohomology on the quotient stack which includes the unstable locus~\cite{MR1792291}.
To do so, we~recall in \cref{section:moduli}
how~$\moduli^{\theta\stable}(Q,\tuple{d})$ is constructed as a quotient
of the Zariski-open~$\repspace[\theta\stable]{Q,\tuple{d}}\subset\repspace{Q,\tuple{d}}$,
and how the bundles~$\mathcal{U}_i^\vee\otimes\mathcal{U}_j$ are the descent
of bundles which are defined on the entire~$\repspace{Q,\tuple{d}}$.
Thus we are in a position to check the conditions for Teleman quantization.

These conditions are checked in \cref{section:teleman-for-quiver-moduli}.
We determine the width of the windows in \cref{subsection:width},
and we determine the weights of the bundles~$\mathcal{U}_i^\vee\otimes\mathcal{U}_j$ in \cref{subsection:weights-endomorphism-bundle}.
The proof of \cref{theorem:cohomology-vanishing}
thus reduces to finding appropriate conditions on~$Q$, $\tuple{d}$ and~$\theta$
for which inequality~\eqref{equation:teleman-inequality-for-real} in the statement of Teleman quantization holds,
which is done in \cref{subsection:cohomology-vanishing}.
We illustrate this in \cref{example:kronecker-moduli-space}.

Four remarks are in order.

\begin{remark}
A crucial intermediate result is \cref{corollary:determinant-weight},
which determines the width of the windows in the setup for Teleman quantization.
Whilst this paper focuses on the weights of~$\mathcal{U}^\vee_i \otimes \mathcal{U}_j$,
there are other bundles to consider from this perspective,
e.g., as~in \cite{MR4894756}.
\end{remark}

\begin{remark}
\label{remark:strong-ample-stability}
The condition that~$\tuple{d}$ is strongly amply~$\theta$\dash stable
in \cref{theorem:cohomology-vanishing} can in fact be replaced by
the (weaker) condition
that the inequality \eqref{equation:teleman-inequality-for-real}
holds for the $1$-parameter subgroups indexed by the Harder--Narasimhan strata.
The reason for opting for strong ample stability in the statement
is that it is easier to verify in practice,
and not directly tied to the problem under consideration. In~\cref{subsection:strong-ample-stability} we describe many situations
in which strong ample stability holds.
An explicit example where strong ample stability does not hold,
but the (weaker) condition from \eqref{equation:teleman-inequality-for-real} does,
is given in \cref{example:strict-ample-stability-rigid}.
\end{remark}

\begin{remark}
\label{remark:ample-stability}
By \cref{lemma:teleman-implies-ample-stability}
we have that \eqref{equation:teleman-inequality-for-real}
implies that~$\tuple{d}$ is amply~$\theta$\dash stable,
i.e., the unstable locus has codimension at least~2.
We have not found a single example which is amply~$\theta$-stable
and for which \eqref{equation:teleman-inequality-for-real} does not hold.

Optimistically, the conclusion of \cref{theorem:cohomology-vanishing} holds
without any condition like \eqref{equation:teleman-inequality-for-real}:
the existence of a universal representation (required for the statement to make sense)
should be enough.
However, as~explained by \cref{lemma:teleman-implies-ample-stability},
one will need new tools beyond Teleman quantization for its proof in this generality.
\end{remark}

\begin{remark}
In \cref{example:twisted-schofield}
we explain how one might be able to consider \cref{theorem:cohomology-vanishing}
in the absence of a universal representation,
replacing it by a twisted universal representation
which is defined on the stable locus.
\end{remark}

We will now discuss three applications of this cohomology vanishing.

\subsubsection*{Schofield's conjecture}
The following conjecture is attributed to Schofield~\cite[p.\,80]{MR1428456}.
\begin{alphaconjecture}[Schofield]
\label{conjecture:schofield}
Let~$Q$ be an acyclic quiver,
and~$\tuple{d}$ a dimension vector.
Let~$\theta$ be a stability parameter
such that~$\modulispace[\theta\stable]{Q,\tuple{d}}$
is a smooth projective variety.
Then the universal representation~$\mathcal{U}$
is a partial tilting bundle on~$\modulispace[\theta\stable]{Q,\tuple{d}}$,
i.e.,
\begin{equation*}
\Ext_{\modulispace[\theta\stable]{Q,\tuple{d}}}^{\geq 1}(\mathcal{U},\mathcal{U})=0.
\end{equation*}
\end{alphaconjecture}
There is a second part to the conjecture,
which states that
this partial tilting bundle can be completed to a tilting bundle.
We will not address the second part in this paper.

The condition that~$Q$ is acyclic in the statement of Schofield's conjecture
is actually superfluous,
and was presumably included to ensure that the moduli space is smooth and projective.
Then the two parts of the conjecture can be seen as a generalization
of Beilinson's description of the derived category of~$\mathbb{P}^n$.

Using the isomorphism
\begin{equation}
\label{equation:Ext-as-cohomology}
\Ext_{\modulispace[\theta\stable]{Q,\tuple{d}}}^{\geq 1}(\mathcal{U},\mathcal{U})
\cong\bigoplus_{i,j\in Q_0}\HH^{\geq 1}(\modulispace[\theta\stable]{Q,\tuple{d}},\mathcal{U}_i^\vee\otimes\mathcal{U}_j),
\end{equation}
\cref{theorem:cohomology-vanishing}
is thus equivalent to Schofield's conjecture.
We encode this reformulation as the following corollary.

\begin{alphacorollary}
\label{corollary:schofield}
Let~$Q$ be a quiver,
let~$\tuple{d}$ be a dimension vector,
and let~$\theta$ be a stability parameter
such that~$\tuple{d}$ is~$\theta$\dash coprime.
Assume furthermore that~$\tuple{d}$ is strongly amply~$\theta$\dash stable.
Then Schofield's conjecture holds,
i.e.,~$\mathcal{U}$ is a partial tilting bundle.
\end{alphacorollary}

Let us now assume that~$Q$ is acyclic, so~that we are only dealing with smooth projective varieties
and finite-dimensional algebras.
\Cref{corollary:schofield} makes it an interesting question to
compute the endomorphism algebra~$\End_{\modulispace[\theta\stable]{Q,\tuple{d}}}(\mathcal{U},\mathcal{U})$,
and thus understand whether the functor
\begin{equation}
\label{equation:fully-faithful-functor}
\Phi_{\mathcal{U}}\colon
\derived^\bounded(\field Q)\to\derived^\bounded(\modulispace[\theta\stable]{Q,\tuple{d}}):
V\mto\RsheafHom_{\mathcal{O}_XQ}(\mathcal{U},-\otimes\mathcal{O}_{\modulispace[\theta\stable]{Q,\tuple{d}}})
\end{equation}
given by the universal representation is fully faithful,
i.e.,
whether
\begin{equation}
\label{equation:End-is-kQ}
\End_{\modulispace[\theta\stable]{Q,\tuple{d}}}(\mathcal{U},\mathcal{U})
\cong
\field Q
\end{equation}
by the morphism naturally induced by~$\Phi_{\mathcal{U}}$. In~the thin case,
i.e.,
when~$\tuple{d}=(1,\ldots,1)$,
and with respect to the canonical stability condition,
the isomorphism \eqref{equation:End-is-kQ} was established by Altmann--Hille in~\cite[Th.\,1.3]{MR1688469}.
Another known case is that of quiver flag varieties,
which are quiver moduli for acyclic quivers with a unique source~$i_0$
and dimension vector~$\tuple{d}$ such that~$d_{i_0}=1$. In~\cite[Th.\,1.2]{MR2772068} Craw shows that both parts of Schofield's conjecture hold,
i.e.,~$\mathcal{U}$ is partial tilting and can be extended to a tilting bundle,
whilst Craw--Ito--Karmazyn show in \cite[Ex.\,2.9]{MR3803802}
that \eqref{equation:End-is-kQ} holds in this setting.

The isomorphism \eqref{equation:End-is-kQ}
beyond the thin and canonical case
is studied by a subset of the authors in~\cite{MR4883104},
building upon the results of this paper.
The main result in op.~cit.,
which uses methods very different from the ones used in this paper,
is that,
under the appropriate conditions,
there exists explicit isomorphisms
\begin{equation*}
\Hom_{\modulispace[\theta\stable]{Q,\tuple{d}}}(\mathcal{U}_i,\mathcal{U}_j)
\cong
\HH^0(\modulispace[\theta\stable]{Q,\tuple{d}},\mathcal{U}_i^\vee\otimes\mathcal{U}_j)
\cong
e_j\field Qe_i,
\end{equation*}
for all~$i,j\in Q_0$.
This result thus naturally complements \cref{theorem:cohomology-vanishing}.

In fact, the fully-faithfulness in~\eqref{equation:fully-faithful-functor}
has a counterpart for moduli of vector bundles on curves.
There exists a rich dictionary between
results for moduli spaces of quiver representations,
and for moduli spaces of vector bundles on curves,
see, e.g.,~\cite{MR3882963} for a survey on some aspects of this dictionary.
Taking this dictionary for granted,
one can take a result for one type of moduli spaces
and try to obtain the analogous result for the other type.

For moduli spaces of vector bundles
the analogue of~\eqref{equation:fully-faithful-functor}
is the Fourier--Mukai functor
\begin{equation*}
\Phi_{\mathcal{E}}\colon\derived^\bounded(C)\to\derived^\bounded(\mathrm{M}_C(r,\mathcal{L})),
\end{equation*}
where~$\mathcal{E}$ is the universal vector bundle on~$C\times\mathrm{M}_C(r,\mathcal{L})$,
and~$\mathrm{M}_C(r,\mathcal{L})$ is the moduli space of
stable vector bundles of rank~$r$ and determinant~$\mathcal{L}$
such that~$\gcd(r,\deg\mathcal{L})=1$
and~$C$ is a smooth projective curve of genus~$g\geq 2$.
Its fully-faithfulness is shown in various levels of generality in~\cite{MR3764066,MR3954042,MR3713871,MR4651618}.

\subsubsection*{Rigidity}
The next application is inspired by the same curve-quiver dictionary.
Recall that in the seminal paper~\cite{MR0384797} Narasimhan--Ramanan show
that the deformation theory of the curve
is the same as that of the moduli space~$\mathrm{M}_C(r,\mathcal{L})$,
again under the assumption that~$\gcd(r,\deg\mathcal{L})=1$ to ensure the moduli space is smooth and projective.
More precisely,
they establish
\begin{equation}
\label{equation:narasimhan-ramanan}
\HH^i(\mathrm{M}_C(r,\mathcal{L}),\tangent_{\mathrm{M}_C(r,\mathcal{L})})
\cong
\HH^i(C,\tangent_C)
\end{equation}
which is well-known to be~$(3g-3)$-dimensional for~$i=1$,
and vanishes for~$i\neq 1$.
For quiver representations and their moduli,
the deformation theory of the quiver
is to be interpreted as the deformation theory of the path algebra of the quiver,
and thus we are interested in its (second) Hochschild cohomology.
By~\cite[\S1.6]{MR1035222} we have that
\begin{equation}
\label{equation:happel}
\hochschild^2(\field Q)=0,
\end{equation}
if~$Q$ is acyclic,
thus the path algebra is \emph{rigid} as an associative algebra.
Under our dictionary between quivers and curves
the analogue of~\eqref{equation:narasimhan-ramanan}
becomes an isomorphism
\begin{equation*}
\HH^1(\modulispace[\theta\stable]{Q,\tuple{d}},\tangent_{\modulispace[\theta\stable]{Q,\tuple{d}}})
\cong
\hochschild^2(\field Q),
\end{equation*}
and thus by the vanishing in~\eqref{equation:happel}
that~$\modulispace[\theta\stable]{Q,\tuple{d}}$
is \emph{rigid as a variety}.
This is proved in \cref{subsection:rigidity}.
The global sections of~$\tangent_{\modulispace[\theta\stable]{Q,\tuple{d}}}$
should be related to~$\hochschild^1(\field Q)$,
which is studied in~\cite{MR4883104}.

Because~\eqref{equation:happel} actually reads~$\hochschild^{\geq 2}(\field Q)$
we in fact expect higher cohomology vanishing for the tangent bundle.
We can prove the rigidity and further vanishing using \cref{theorem:cohomology-vanishing}
and the 4-term sequence~\eqref{equation:4-term-exact-sequence}.

\begin{alphacorollary}
\label{corollary:rigidity}
Let~$Q$ be an acyclic quiver,
let~$\tuple{d}$ be a dimension vector,
and let~$\theta$ be a stability parameter,
such that~$\tuple{d}$ is~$\theta$-coprime.
Assume furthermore that $\tuple{d}$ is strongly amply~$\theta$-stable.
Then~$\modulispace[\theta\stable]{Q,\tuple{d}}$ satisfies
\begin{equation*}
\HH^{\geq 1}(\modulispace[\theta\stable]{Q,\tuple{d}},\tangent_{\modulispace[\theta\stable]{Q,\tuple{d}}})=0,
\end{equation*}
so in particular it is (infinitesimally) rigid,
i.e.,~$\HH^1(\modulispace[\theta\stable]{Q,\tuple{d}},\tangent_{\modulispace[\theta\stable]{Q,\tuple{d}}})=0$.
\end{alphacorollary}

Thus, it is not possible to deform a quiver moduli space,
or put more colloquially: ``quiver moduli have no moduli''.
The only cases where this rigidity was known
were situations in which an explicit description of the moduli space was known
(e.g., for certain quiver moduli attached to Kronecker or subspace quivers),
or in the case of Fano toric quiver moduli, by~applying Danilov–Steenbrink–Bott vanishing,
see, e.g.,~\cite[Th.\,2.4(i)]{MR1916637}.

\subsubsection*{Finiteness}
In~\cite{2303.08522v2} Domokos proves a result which interacts
in an interesting way with \cref{corollary:rigidity}. In~op.~cit.~it is shown that,
if one fixes the dimension,
there exist only \emph{finitely} many isomorphism classes of quiver moduli of that dimension,
where a priori the number could also be countable
(as the choice of quiver, dimension vector, and stability chamber is
countable, \resp countable, \resp finite).
There are three variations on the finiteness statement in op.~cit.,
with different conditions on the dimension vectors that one considers.

One can also \emph{reverse} the flow of information:
by the boundedness of Fano varieties \cite[Th.\,0.2]{MR1189503}
we know that in a given dimension there are only finitely many \emph{rigid} Fano varieties.
Applying \cref{corollary:rigidity} to Fano quiver moduli (for strongly amply stable dimension vectors)
one obtains that there are only finitely many such quiver moduli in a given dimension,
giving a fourth variation on the finiteness statement from \cite{2303.08522v2}.

\smallskip

The computation of the global sections of the tangent bundle,
i.e., the vector fields of the moduli space,
is addressed in~\cite{MR4883104}. It~is shown that,
under favorable circumstances,
there exists an isomorphism~$\hochschild^1(\field Q)\cong\HH^0(\modulispace[\theta\stable]{Q,\tuple{d}},\tangent_{\modulispace[\theta\stable]{Q,\tuple{d}}})$, as~predicted by the dictionary between quivers and curves.

\subsubsection*{Height-zero moduli spaces}
The final application involves moduli spaces of sheaves on~$\mathbb{P}^2$,
and in particular those of height zero, as~introduced by Drezet~\cite{MR0916199,MR0944602,MR0915184}.
These are moduli spaces~$\moduli_{\mathbb{P}^2}(r,\mathrm{c}_1,\mathrm{c}_2)$
of (semi)stable sheaves on~$\mathbb{P}^2$
which can be characterized by having~$\Pic\moduli_{\mathbb{P}^2}(r,\mathrm{c}_1,\mathrm{c}_2)\cong\mathbb{Z}$~\cite[Th.\,2]{MR0915184}.
The precise definition is recalled in \cref{subsection:height-zero}.

Using the parameters~$(r,\mathrm{c}_1,\mathrm{c}_2)=(1,0,n)$
we describe the Hilbert scheme of~$n$ points on~$\mathbb{P}^2$
as a fine moduli space of stable sheaves.
In~\cite[Th.\,1.2]{MR3397451} it is shown that the Fourier--Mukai functor~$\Phi_{\mathcal{I}}\colon\derived^\bounded(\mathbb{P}^2)\to\derived^\bounded(\Hilb^n\mathbb{P}^2)$,
where~$\mathcal{I}$ is the universal ideal sheaf,
is fully faithful for all~$n\geq 2$.
In~\cite[Prop.\,29]{MR3950704} it is shown that this implies that~$\HH^1(\Hilb^n\mathbb{P}^2,\tangent_{\Hilb^n\mathbb{P}^2})\cong\hochschild^2(\mathbb{P}^2)$ is non-zero.
As mentioned in Remark~30 of op.~cit.,
the same method to compute the infinitesimal deformations
works verbatim for other smooth projective moduli spaces,
provided the Fourier--Mukai functor is fully faithful.

This gives us examples where the functor cannot be fully faithful,
using the correspondence between moduli spaces of height zero
and certain Kronecker moduli~\cite[Th.\,2]{MR0916199}.
We prove the following in \cref{subsection:height-zero}.

\begin{alphacorollary}
\label{corollary:height-zero}
Let~$\moduli_{\mathbb{P}^2}(r,\mathrm{c}_1,\mathrm{c}_2)$ be a smooth projective
fine moduli space of stable sheaves on~$\mathbb{P}^2$,
with universal object~$\mathcal{E}$.
If it is of height zero,
then
\begin{equation}
\label{equation:height-zero-functor}
\Phi_{\mathcal{E}}\colon\derived^\bounded(\mathbb{P}^2)\to\derived^\bounded(\moduli_{\mathbb{P}^2}(r,\mathrm{c}_1,\mathrm{c}_2))
\end{equation}
is not fully faithful.
\end{alphacorollary}

The special case of the corollary where~$(r,\mathrm{c}_1,\mathrm{c}_2)=(4,1,3)$
was established in~\cite{MR4223526}.
It is also possible to obtain the result in \cref{corollary:height-zero} using other methods, as~explained in \cref{remark:associated-exceptional-bundle},
with ingredients from the original papers by Drezet.

\subsubsection*{Acknowledgements}

The first author would like to thank Ben Gould, Dmitrii Pedchenko and Fabian Reede
for conversations about moduli spaces of height zero
related to \cref{remark:associated-exceptional-bundle}.
We want to thank Alastair Craw for interesting discussions regarding the case of quiver flag varieties.
We want to thank Daniel Halpern-Leistner and Sebastían Torres for interesting discussions
regarding Teleman quantization.

\section{Moduli spaces of quiver representations}
\label{section:moduli}

We will first set up the notation,
and recall the construction of various moduli spaces of quiver representations
with their universal objects.
Throughout we will let~$\field$ be an algebraically closed field of characteristic~0.
Initially this is not necessary, but it will be for the main results.

A \emph{quiver} $Q$ is a finite directed graph,
with set of vertices $Q_0$ and set of arrows $Q_1$,
together with two maps $\source, \target\colon Q_1 \to Q_0$,
which assign to an arrow $a \in Q_1$ its source and target, respectively.

A \emph{representation} $M$ of $Q$ over a field $\field$
is the data of a finite-dimensional vector space~$M_i$
for each vertex~$i \in Q_0$
together with a linear map~$M_a\colon M_{\source(a)} \to M_{\target(a)}$
for each arrow~$a \in Q_1$.
For a representation~$M$ of~$Q$, we~define the \emph{dimension vector} $\dimvect M \in \smash{\mathbb{N}^{Q_0}}$ as the tuple $(\dim M_i)_{i \in Q_0}$.

We will briefly recall the construction of the moduli stacks and spaces,
working over the field~$\field$.
For a more formal definition,
valid over arbitrary bases,
the interested reader can consult~\cite[\S3]{2210.00033v1}.
Given a fixed dimension vector $\tuple{d} \in \smash{\mathbb{N}_0^{Q_0}}$,
where we write~$\tuple{d}=(d_i)_{i\in Q_0}$, we~define the \emph{representation variety}
as the affine space
\begin{equation}
\label{equation:representation-variety}
\repspace{Q,\tuple{d}}\colonequals \prod_{a \in Q_1} \Mat_{d_{\target(a)} \times d_{\source(a)}}(\field).
\end{equation}
A point of $\repspace{Q,\tuple{d}}$ is the same as a representation of $Q$ on the collection of vector spaces $\field^{d_i}$.
Let the reductive linear algebraic group $\group=\group(\field)$ be defined by
\begin{equation*}
\group(\field)\colonequals \prod_{i \in Q_0} \GL_{d_i}(\field).
\end{equation*}
This group acts on the left on $\repspace{Q,\tuple{d}}$ via change of basis;
more precisely,
$g = (g_i)_{i \in Q_0}\in\group$ acts on the point~$M = (M_a)_{a \in Q_1}$ by
\begin{equation*}
g\cdot M \colonequals (g_{\target(a)}M_a\,g_{\source(a)}^{-1})_{a \in Q_1}.
\end{equation*}
Two elements of $\repspace{Q,\tuple{d}}$ lie in the same $\group$-orbit
if and only if the corresponding representations of $Q$ are isomorphic.
Consider the closed central subgroup \hbox{$\Delta = \{(z\cdot \mathrm{I}_{d_i})_{i \in Q_0} \mid z \in \Gm\}$} of $\group$.
Its elements act trivially on $\repspace{Q,\tuple{d}}$,
whence the action on $\repspace{Q,\tuple{d}}$ descends to an action of the reductive group
\begin{equation*}
\PG_{\tuple{d}} \colonequals \group/\Delta.
\end{equation*}

For later use we also recall the Euler pairing of the quiver,
which for~$\alpha, \beta \in \mathbb{Z}^{Q_0}$
is defined as
\begin{equation*}
\langle \alpha, \beta \rangle = \sum_{i \in Q_0} \alpha_i \beta_i ~- \sum_{a \in Q_1} \alpha_{\source(a)} \beta_{\target(a)}.
\end{equation*}
Note that the notation does not include $Q$, as~it will be clear from the context.

\subsubsection*{Moduli of (semi)stable representations}
Let $\theta\in\Hom(\mathbb{Z}^{Q_0},\mathbb{Z})$ be a \emph{stability parameter}.
We will identify $\Hom(\mathbb{Z}^{Q_0},\mathbb{Z})$ with its dual~$\mathbb{Z}^{Q_0}$ by the dot product.
Assume that $\theta(\tuple{d}) = 0$.
We define a notion of stability with respect to $\theta$,
following King~\cite{MR1315461} (up to a different sign convention).
\begin{definition}
\label{definition:semistability}
A representation $M$ of $Q$ of dimension vector $\tuple{d}$
is called $\theta$\nobreakdash-\emph{semistable} (\resp $\theta$\dash\emph{stable})
if any proper non-zero subrepresentation $N$ of $M$
satisfies the inequality~$\theta(\dimvect N) \leq 0$ (\resp $\theta(\dimvect N) < 0$).
\end{definition}

Let
\begin{equation*}
\repspace[\theta\stable]{Q,\tuple{d}}
\subseteq\repspace[\theta\semistable]{Q,\tuple{d}}
\subseteq\repspace{Q,\tuple{d}}
\end{equation*}
denote the Zariski open subsets of stable and semistable representations, respectively.
Note that the semistable locus may be empty,
or the stable locus may be empty while the semistable locus is not.

An important choice of stability parameter is
the \emph{canonical stability parameter}~$\theta_\can$ for a dimension vector~$\tuple{d}$,
which is given by the morphism~$\langle\tuple{d},-\rangle-\langle-,\tuple{d}\rangle$,
i.e.,
the partial evaluation of the antisymmetrized Euler form~$\{-,-\}$.

We associate the character $\chi_\theta$ of $\group$
to the stability parameter~$\theta$, by~defining
\begin{equation*}
\chi_\theta(g) = \prod_{i \in Q_0} \det(g_i)^{-\theta_i};
\end{equation*}
note the minus sign in the exponent.
As $\theta(\tuple{d}) = 0$,
the character $\chi_{\theta}$ vanishes on~$\Delta$, so~it is also a character of $\PG_{\tuple{d}}$.
Let $L(\theta)$ be the trivial line bundle on $\repspace{Q,\tuple{d}}$
equipped with the $\PG_\tuple{d}$-linearization given by $\chi_\theta$.
King shows in~\cite[Th.\,4.1]{MR1315461} that $\theta$-semistability
agrees with semistability with respect to $L(\theta)$,
and $\theta$-stability is the same as proper stability in the sense of Mumford's GIT. In~op.~cit., the action by~$\group$ is used, so~the points are not properly stable in the sense of Mumford,
but the only difference is a common~$\Gm$-stabilizer.

\begin{definition}
\label{definition:moduli}
We define the $\theta$-semistable and $\theta$-stable \emph{moduli spaces} as
\begin{align*}
\modulispace[\theta\semistable]{Q,\tuple{d}} & \colonequals \repspace{Q,\tuple{d}}/\!\!/_{L(\theta)} \PG_{\tuple{d}}, \\
\modulispace[\theta\stable]{Q,\tuple{d}} & \colonequals \repspace{Q,\tuple{d}}/_{L(\theta)} \PG_{\tuple{d}}.
\end{align*}
Moreover, we~define the $\theta$-semistable and $\theta$-stable \emph{moduli stacks} as
\begin{align*}
\modulistack[\theta\semistable]{Q,\tuple{d}} & \colonequals [\repspace[\theta\semistable]{Q,\tuple{d}}/\group], \\
\modulistack[\theta\stable]{Q,\tuple{d}} & \colonequals [\repspace[\theta\stable]{Q,\tuple{d}}/\group].
\end{align*}
\end{definition}

In the case $\theta = 0$, we~write $\modulispace[\semisimple]{Q,\tuple{d}}$
for the moduli space of $0$-semistable representations, as~well as $\modulistack{Q,\tuple{d}} = \modulistack[0\semistable]{Q,\tuple{d}}$
for the moduli stack.
These parameterize the semisimple representations of dimension vector~$\tuple{d}$.

The relations between the affine, semistable and stable quotients are summarized in the following diagram:
\begin{equation*}
\begin{tikzcd}[row sep=small, ampersand replacement=\&]
\repspace[\theta\stable]{Q,\tuple{d}} \arrow[r,hook] \arrow[d] \& \repspace[\theta\semistable]{Q,\tuple{d}} \arrow[r,hook] \arrow[d] \& \repspace{Q,\tuple{d}} \arrow[d] \\
\modulistack[\theta\stable]{Q,\tuple{d}} \arrow[r,hook] \arrow[d] \& \modulistack[\theta\semistable]{Q,\tuple{d}} \arrow[r,hook] \arrow[d] \& \modulistack{Q,\tuple{d}} \arrow[d] \\
\modulispace[\theta\stable]{Q,\tuple{d}} \arrow[r, hook] \& \modulispace[\theta\semistable]{Q,\tuple{d}} \arrow[r] \& \modulispace[\semisimple]{Q,\tuple{d}}.
\end{tikzcd}
\end{equation*}
The arrows on the first row are open immersions.
So are the leftmost arrows in the second and third row.
The morphism $\modulistack[\theta\stable]{Q,\tuple{d}} \to \modulispace[\theta\stable]{Q,\tuple{d}}$
is a $\Gm$-gerbe,
the composition $\repspace[\theta\stable]{Q,\tuple{d}} \to \modulispace[\theta\stable]{Q,\tuple{d}}$
is a principal $\PG_\tuple{d}$-bundle in the fpqc topology,
and $\modulispace[\theta\semistable]{Q,\tuple{d}} \to \modulispace[\semisimple]{Q,\tuple{d}}$ is projective.
The semisimple moduli space $\modulispace[\semisimple]{Q,\tuple{d}}$ is affine,
and if~$Q$ is acyclic
this moduli space is isomorphic to~$\Spec\field$
by the classification of simple modules.
The moduli space $\modulispace[\theta\stable]{Q,\tuple{d}}$ is smooth.

\begin{definition}
A dimension vector $\tuple{d} \in \mathbb{N}^{Q_0}$ is called \emph{indivisible}
if $\gcd_{i \in Q_0}(d_i) = 1$.
Furthermore, it is called \emph{$\theta$-coprime} if,
for any $\tuple{d'} \in \mathbb{N}^{Q_0}$ such that $\tuple{d'} \leq \tuple{d}$, we~have $\theta(\tuple{d'}) \neq 0$,
unless $\tuple{d'} \in \{0,\tuple{d}\}$.
\end{definition}
We see that if $\tuple{d}$ is $\theta$-coprime, then it is indivisible.
Moreover, if $\tuple{d}$ is $\theta$-coprime,
then every $\theta$-semistable representation of dimension vector $\tuple{d}$ is $\theta$-stable. In~particular, if the quiver $Q$ is acyclic and if~$\tuple{d}$ is~$\theta$\dash coprime,
the moduli space~$\modulispace[\theta\stable]{Q,\tuple{d}}$ is a smooth projective variety.
If moreover~$\theta=\theta_\can$ (or more generally, $\theta$ is in the same chamber as~$\theta_\can$)
then~$\modulispace[\theta\stable]{Q,\tuple{d}}$ is even a Fano variety \cite{MR4352662},
thus explaining why this particular choice is especially interesting.

We also recall the following definition from~\cite[Def.\,4.1]{MR3683503}.
\begin{definition}
\label{definition:ample-stability}
A dimension vector~$\tuple{d}$ is said to be \emph{amply~$\theta$-stable} if
\begin{equation*}
\operatorname{codim}_{\Rep(Q,\tuple{d})}(\Rep(Q,\tuple{d})\setminus\Rep^{\theta\stable}(Q,\tuple{d}))\geq 2.
\end{equation*}
\end{definition}
This guarantees that the Picard rank of~$\modulispace[\theta\stable]{Q,\tuple{d}}$ will be maximal,
i.e., it is equal to~$\#Q_0-1$.
We will introduce a stronger version of this condition in \cref{definition:strong-ample-stability}.

\subsubsection*{Universal representations}
For every $i \in Q_0$,
let $U_i$ be the trivial vector bundle of rank~$d_i$ on $\repspace{Q,\tuple{d}}$.
We equip it with an action of $\group$.
For a vector $v \in \field^{d_i} = (U_i)_M$ in the fiber over a point $M \in \repspace{Q,\tuple{d}}$,
the group element $g=(g_i)_{i\in Q_0} \in \group$ acts by
\begin{equation*}
g \cdot v = g_iv
\end{equation*}
which lies in the fiber over $g\cdot M$.
For an arrow $a \in Q_1$,
let $U_a\colon U_{\source(a)} \to U_{\target(a)}$ be
the morphism which on the fibers over a point $M \in \repspace{Q,\tuple{d}}$
sends a vector $v \in (U_{\source(a)})_M$ to $M_a(v) \in (U_{\target(a)})_M$.
This morphism is clearly~$\group$\dash equivariant.

The bundles $U_i$ do not descend to $\modulispace[\theta\stable]{Q,\tuple{d}}$
because the stabilizer $\Delta$ does not act trivially on the fibers.
We need to twist by certain line bundles to achieve this.

Let~$\tuple{a} \in \mathbb{Z}^{Q_0}$.
Define~$L(\tuple{a})$ as the trivial line bundle on~$\Rep(Q,\tuple{d})$, so~that on its fibers the element~$g \in \group$ acts by
the character~$\chi_\tuple{a}(g) = \prod_{i \in Q_0} \det(g_i)^{-a_i}$.
If we assume that~$\tuple{d}$ is indivisible,
then there exists a tuple of integers~$\tuple{a} \in \mathbb{Z}^{Q_0}$
for which~$\tuple{a}\cdot\tuple{d}=\sum_{i \in Q_0}a_i d_i = 1$.

Let $U_i(\tuple{a}) \colonequals U_i \otimes L(\tuple{a})$ on~$\Rep(Q,\tuple{d})$
for~$\tuple{a}$ such that~$\tuple{a}\cdot\tuple{d}=1$. By~the choice of $\tuple{a}$, the stabilizer $\Delta$ acts trivially, so~these vector bundles, once restricted to the $\theta$-stable locus,
descend to the quotient~$\modulispace[\theta\stable]{Q,\tuple{d}}$.
The morphisms~$U_a \otimes \operatorname{id}_{L(\tuple{a})}\colon U_{\source(a)}(\tuple{a}) \to U_{\target(a)}(\tuple{a})$
descend to morphisms~$\mathcal{U}_a(\tuple{a})\colon \mathcal{U}_{\source(a)}(\tuple{a}) \to \mathcal{U}_{\target(a)}(\tuple{a})$.

\begin{definition}
The data $\mathcal{U}(\tuple{a}) = ((\mathcal{U}_i(\tuple{a}))_{i \in Q_0},(\mathcal{U}_a(\tuple{a}))_{a \in Q_1})$
is a representation of $Q$ in vector bundles on $\modulispace[\theta\stable]{Q,\tuple{d}}$. It~is called the \emph{universal representation} (relative to~$\tuple{a}$).
\end{definition}

\begin{remark}
\label{remark:uniqueness-up-to-twist-universal-bundle}
Different choices of $\tuple{a}$ give rise to non-isomorphic universal representations.
However, for~$\tuple{a}$ and~$\tuple{b} \in \mathbb{Z}^{Q_0}$ as before,
i.e., $\tuple{a}\cdot\tuple{d}=\tuple{b}\cdot\tuple{d}=1$, we~have universal representations~$\mathcal{U}(\tuple{a})$ and~$\mathcal{U}(\tuple{b})$.
The line bundle $L(\tuple{b} - \tuple{a})$
descends to a line bundle $\mathcal{L}(\tuple{b} - \tuple{a})$ on $\modulispace[\theta\stable]{Q,\tuple{d}}$,
because~$(\tuple{b}-\tuple{a})\cdot\tuple{d}=0$,
making the stabilizer act trivially,
and it gives rise to the isomorphism $\mathcal{U}(\tuple{a}) \otimes \mathcal{L}(\tuple{b} - \tuple{a}) \cong \mathcal{U}(\tuple{b})$.
\end{remark}

\section{Teleman quantization for quiver moduli}
\label{section:teleman-for-quiver-moduli}
The following section is the technical heart of the paper.
We will explain how to set up Teleman quantization for quiver moduli, by~recalling the Hesselink stratification
and relating it to the Harder--Narasimhan stratification.
This allows us to compute the width of the windows in \cref{subsection:width},
and the weights of the endomorphisms of the universal representation in \cref{subsection:weights-endomorphism-bundle}.

\subsection{The Hesselink stratification and Teleman's quantization theorem}
\label{subsection:hesselink-and-teleman}
Let $G$ be a linearly reductive algebraic group,
and let $R$ be an affine variety over $\field$ on which~$G$ acts. In~our application to quiver moduli we will let~$G$ be~$\group$ or~$\PG_{\tuple{d}}$,
and~$R$ the representation variety~\eqref{equation:representation-variety}.

Let~$\lambda\colon\Gm\to G$ be a $1$-parameter subgroup,
and let~$\chi\colon G \to \Gm$ be a character of~$G$.
We denote by~$\langle\chi,\lambda\rangle$ the integer exponent in the identity~$\chi \circ \lambda (z) = z^{\langle\chi,\lambda\rangle}$.
Recall that the Hilbert--Mumford criterion for semistability states the following.

\begin{theorem}
\label{theorem:hilbert-mumford}
A point $x \in R$ is $\chi$-semistable if and only if $ \langle\chi,\lambda\rangle \geq 0$
for every $1$-parameter subgroup~$\lambda\colon \Gm \to G$
for which $\lim_{z \to 0} \lambda(z)x$ exists.
\end{theorem}
Given an unstable point $x \in R$, Kempf finds in~\cite{MR0506989} a $1$-parameter subgroup which is
``most responsible'' for its instability.
This construction was used by Hesselink in~\cite{MR0514673}
to obtain a stratification of $R$ into locally closed subsets.
The Hesselink stratification is also the basis for Teleman's quantization theorem.
We will thus briefly recall this theory in the following.
A comprehensive reference is~\cite[\S 12 \& \S 13]{MR0766741}.

With $\cocharacters(G)$ we denote the set of $1$-parameter subgroups of $G$. It~is acted upon by~$G$ by conjugation.
Fix a maximal torus $T \subseteq G$ and let $W$ be the corresponding Weyl group.
Then $\cocharacters(T)$ is a free abelian group. The Weyl group $W$ acts on $\cocharacters(T)$ and we have a bijection
\begin{equation}\label{equation:cocharacters}
\cocharacters(T)/W \cong \cocharacters(G)/G.
\end{equation}
Let $(-,-)$ be a $W$-invariant inner product on the vector space $\cocharacters(T)_\mathbb{R}$
such that $(\lambda,\lambda) \in \mathbb{Z}$ for all $1$-parameter subgroups $\lambda$ of $T$.
The induced norm $\|{-}\|$ is also $W$\nobreakdash-invariant.
Therefore, we~may extend the norm to $1$-parameter subgroups $\lambda$ of $G$
by setting $\|\lambda\|\colonequals\|g\lambda g^{-1}\|$ where $g \in G$ is such that $g\lambda g^{-1}$ lies in $T$.

\begin{definition}
Let $x \mkern-4mu\in\mkern-5mu R$. We define the \emph{normalized Hilbert--Mumford weight}~of~$x$~as
\begin{equation*}
m(x)\colonequals \inf\Bigl\{ \frac{\langle\chi,\lambda\rangle}{\|\lambda\|}\Bigm| 1 \neq \lambda \in \cocharacters(G) \text{ such that } \lim_{z\to 0} \lambda(z)x \text{ exists}\Bigr\},
\end{equation*}
and the set of $1$-parameter subgroups associated to $x$ as
\begin{equation*}
\Lambda(x) \colonequals \Bigl\{ \lambda \in \cocharacters(G)\Bigm|\lambda \text{ is primitive, } \lim_{z\to 0} \lambda(z)x \text{ exists, and } \frac{\langle\chi,\lambda\rangle}{\|\lambda\|} = m(x) \Bigr\}.
\end{equation*}
\end{definition}

For a $1$-parameter subgroup $\lambda$ of $G$ we define
\begin{align}%\label{definition:P_lambda}
L_\lambda &:= \{ g \in G\mid \lambda(z)g\lambda(z)^{-1} = g \text{ for all } z \in \Gm \}, \text{ and}\\
P_\lambda &:= \{ g \in G\mid \lim_{z \to 0} \lambda(z)g\lambda(z)^{-1} \text{ exists}\}.
\end{align}
These are closed subgroups of $G$.
Moreover, $P_\lambda$ is a parabolic subgroup with Levi factor $L_\lambda$.

We define two subsets of $R$ as follows,
where the first is the fixed locus of~$\lambda$:
\begin{align*}
R_\lambda & \colonequals \{ x \in R \mid \lambda(z)x = x \text{ for all } z \in \Gm \}, \text{ and} \\
R_\lambda^+ & \colonequals \{ x \in R \mid \lim_{z \to 0}\lambda(z)x \text{ exists}\}.
\end{align*}
Both subsets are closed in the Zariski topology.
The group $L_\lambda$ acts on $R_\lambda$, while $P_\lambda$ acts on $R_\lambda^+$.
The morphism
\begin{equation}
\label{equation:p-lambda}
p_\lambda\colon R_\lambda^+ \to R_\lambda:x\mto\lim_{z \to 0}\lambda(z)x
\end{equation}
is $P_\lambda$-equivariant via $P_\lambda \to L_\lambda$.

\begin{definition}
Let $\lambda$ be a $1$-parameter subgroup of $G$,
and let $[\lambda]$ be its $G$\nobreakdash-conjugacy class inside~$\cocharacters(G)$.
We define the following subsets of $R$:
\begin{enumerate}
\item $S_{[\lambda]}\colonequals \{ x \in R \mid \Lambda(x) \cap [\lambda] \neq \emptyset \}$,
which is called the \emph{Hesselink stratum} of $[\lambda]$;
\item $\blade[\lambda]\colonequals \{ x \in R \mid \lambda \in \Lambda(x) \}\subset S_{[\lambda]}$,
which is called the \emph{blade} of $\lambda$; and
\item $Z_\lambda\colonequals \{ x \in R_\lambda \mid \lambda \in \Lambda(x)\} = \blade[\lambda] \cap R_\lambda\subset\blade[\lambda]$,
which is called the \emph{limit set} of $\lambda$.
\end{enumerate}
\end{definition}

Now we proceed to describe the Hesselink stratification. It~is an algebraic analog of the Kempf--Ness stratification.

\begin{theorem}[Hesselink]
\label{theorem:hesselink}
The unstable locus (or null cone) inside $R$ admits a decomposition
\begin{equation}
\label{equation:hesselink-stratification}
R\setminus R^{\chi\semistable} = \bigsqcup_{[\lambda]} S_{[\lambda]},
\end{equation}
as a finite disjoint union into $G$-invariant locally closed subsets
ranging over all the $G$-conjugacy classes of primitive $\lambda \in \cocharacters(G)$
such that $\langle \chi,\lambda \rangle < 0$.
Moreover,
\begin{enumerate}
\item Each limit set $Z_\lambda$ is $L_\lambda$-invariant and open inside $R_\lambda$;
\item Each blade satisfies $\blade[\lambda] = p_\lambda^{-1}(Z_\lambda)$,
where~$p_\lambda$ is as in~\eqref{equation:p-lambda},
and is hence a $P_\lambda$\nobreakdash-invari\-ant open subset of~$R_\lambda^+$;
\item Each stratum $S_{[\lambda]}$ satisfies $S_{[\lambda]} = G \cdot \blade[\lambda]$; and
\item The action map $\sigma\colon G\times \blade[\lambda] \to G\cdot \blade[\lambda] = S_{[\lambda]}$
induces an isomorphism with the associated fiber bundle,
i.e.,~$G\times^{P_\lambda} \blade[\lambda] \cong S_{[\lambda]}$.
\end{enumerate}
\end{theorem}
The complement of the null cone can be included in the decomposition \eqref{equation:hesselink-stratification}
by using the trivial $1$-parameter subgroup.

\begin{remark}
\label{remark:indexing-set}
By~\eqref{equation:cocharacters},
the disjoint union of~\eqref{equation:hesselink-stratification} may equivalently be taken over
all $W$-conjugacy classes of primitive $1$-parameter subgroups $\lambda$ of $T$
such that $\langle \chi,\lambda \rangle < 0$.
\end{remark}

Now we state the result which is most important for our purposes -- Teleman's quantization theorem.
We do not state its most general version, but only one which will be sufficient for our purposes.
This result was given in great generality by Halpern-Leistner in~\cite[Th.\,3.29]{MR3327537},
extending the original version from \cite[Prop.\,2.11 \& Rem.\,2.12(i)]{MR1792291}.

\begin{theorem}[Teleman]
\label{theorem:teleman-quantization}
Let~$G$ be a reductive algebraic group
acting on a smooth quasiprojective variety~$R$.
Let $\{S_{[\lambda]}\}$ be the Hesselink stratification~\eqref{equation:hesselink-stratification} of the unstable locus
with respect to the character $\chi$ from \cref{theorem:hesselink}.
Assume that all limit sets $Z_\lambda$ are connected.

For each $\lambda$, define $\eta_\lambda$ as the weight of the action of $\lambda$
on the determinant of the conormal bundle restricted to the limit set, i.e.,
\begin{equation}
\label{equation:eta-lambda}
\eta_\lambda\colonequals\weight_\lambda \bigl( (\det \normal_{S_{[\lambda]}/R}^\vee) |_{Z_\lambda} \bigr).
\end{equation}
Let $\mathcal{F}$ be a $G$-linearized vector bundle on $R$.
If, for each $\lambda$, the weights of the action of $\lambda$ on $\mathcal{F}|_{Z_\lambda}$
are strictly less\footnote{Some authors state this result
by taking the limit~$z\to\infty$ rather than~$z\to0$,
which modifies the statement.}
than $\eta_\lambda$,
i.e.,
\begin{equation}
\label{equation:teleman-inequality-general}
\max\left\{ \weight_\lambda\left( \mathcal{F}|_{Z_\lambda} \right) \right\}<\eta_\lambda,
\end{equation}
then the natural map
\begin{equation*}
\HH^k(R,\mathcal{F})^G \to \HH^k(R^{\chi\semistable},\mathcal{F})^G
\end{equation*}
is an isomorphism for all $k \geq 0$.
\end{theorem}

In particular, if~$G$ acts freely on the semistable locus,
then we obtain isomorphisms
\begin{equation*}
\HH^k(R,\mathcal{F})^G \to \HH^k(R{/\!\!/_\chi}\,G,\mathcal{F})
\end{equation*}
for all~$k\geq 0$,
where on the right-hand side~$\mathcal{F}$ denotes the descent to the GIT quotient, as~in \cite[Th.\,3.2]{MR1792291}.

Because the statement in \cref{theorem:teleman-quantization}
is not readily obtained from the statement in \cite[Th.\,3.29]{MR3327537}
we sketch how it can be deduced.
Alternatively,
one can go through the setup of~\cite[Prop.\,2.11]{MR1792291},
taking into account the different sign convention,
and different notation.
\begin{proof}[Sketch of proof of \cref{theorem:teleman-quantization}]
\cite[Th.\,3.29]{MR3327537} gives an isomorphism for the restriction
on a per-stratum basis.
One applies it inductively, as~in the proof of \cite[Th.\,1.1]{MR3327537},
checking the conditions for~$F^\csbullet$ and~$G^\csbullet$
in the notation of op.~cit.~for \emph{all} strata.

To apply this result, we~take~$F^\csbullet=\mathcal{O}_R$,
and~$G^\csbullet=\mathcal{F}$,
with weights~$v=w=0$ for every stratum.
For every stratum we have~$\mathcal{O}_R\in\derived^-([R/G])_{\geq 0}$
as defined in \cite[Def.\,3.24]{MR3327537}
because the weight of the structure sheaf is~0.
Finally, we~have~$\mathcal{F}\in\derived^+([R/G])_{<0}$
by the smoothness assumption on~$R$, so~that we can apply \cite[Lem.\,2.9]{MR3327537},
which translates the inequality
in the definition of~$\derived^+([R/G])_{<0}$
to \eqref{equation:teleman-inequality-general}.
\end{proof}

\begin{remark}
\label{remark:weights}
Note that as~$Z_\lambda$ is supposed to be connected,
and $\Gm$ acts trivially via~$\lambda$ on~$Z_\lambda$,
\cite[Prop.\,2.10]{MR3329192} gives a split exact sequence
\begin{equation*}
0 \to \mathbb{Z} \to \operatorname{Pic}^{\Gm}(Z_\lambda) \to \operatorname{Pic}(Z_\lambda) \to 0.
\end{equation*}
A section is given by equipping a line bundle with the trivial linearization.
The corresponding retraction is the map $\weight_\lambda\colon \operatorname{Pic}^{\Gm}(Z_\lambda) \to \mathbb{Z}$.
This means that a $\Gm$\nobreakdash-lin\-earized line bundle $L$ on $Z_\lambda$
is the same as a line bundle on $Z_\lambda$ together with
a linear action of $\Gm$ on every fiber by the same weight $\weight_\lambda(L)$.
\end{remark}

\Subsection{The Harder--Narasimhan and Hesselink stratifications for quiver moduli}
\label{section:HN-vs-Hesselink}
As in \cref{section:moduli}
we let~$Q$ be a quiver, $\tuple{d}$ a dimension vector
and~$\theta \in\Hom(\mathbb{Z}^{Q_0},\mathbb{Z})$ a stability parameter
such that $\theta(\tuple{d})=0$.
The character~$\chi_\theta$ of~$\group$ defined by~$\chi_\theta(g) = \prod_{i \in Q_0} \det(g_i)^{-\theta_i}$
descends to a character of~$\PG_{\tuple{d}} = \group/\Delta$,
therefore it makes no difference to consider the action of~$\group$ instead of the action of~$\PG_{\tuple{d}}$.

We consider the action of $\PG_{\tuple{d}}$ on the representation space $R = \repspace{Q,\tuple{d}}$
and describe the Hesselink stratification in this case.

For quiver representations, the fifth-named author established in~\cite[Prop.\,3.4]{MR1974891}
the existence of an analogous stratification of $R$, the \emph{Harder--Narasimhan stratification}.
Under the standing assumption that~$\field$ is algebraically closed these two will coincide,
which is the content of \cref{theorem:HN=hesselink}. In~the case of complex numbers,
Hoskins gives a proof in~\cite[Th.\,3.8]{MR3871820}.
Over an arbitrary algebraically closed field,
Zamora gives a proof in~\cite[Th.\,6.3]{MR3199484}.
We recall below how the identification goes,
and explain how all the auxiliary data are related,
because we will need this for later computations.

We make the representation-theoretic setup of \cref{subsection:hesselink-and-teleman}
explicit in our special setting.
Let~$T_{\tuple{d}} \subseteq \group$ be the maximal torus of diagonal matrices.
The lattice $\cocharacters(T_{\tuple{d}})$ of $1$-parameter subgroups of~$T_{\tuple{d}}$
is freely generated by~$\{\lambda_{i,r} \mid i \in Q_0,\ r=1,\ldots,d_i\}$
where~$\lambda_{i,r}(z) \in T_{\tuple{d}}$
is the tuple of matrices where in the $i$th matrix there is a diagonal entry $z$
in the $r$th position.
Let~$(-,-)$ be the inner product on~$\cocharacters(T_{\tuple{d}})_\mathbb{R}$
for which~$\{\lambda_{i,r}\}$ is an orthonormal basis.
This inner product is invariant under the Weyl group~$W = \prod_{i \in Q_0}\mathrm{S}_{d_i}$.
Let $\|{-}\|$ be the corresponding norm.
Note that,
for each $i \in Q_0$,
the restriction of this norm on $\cocharacters(T_{d_i})$ is the standard Euclidean norm.

We also have to consider the maximal torus~$T = T_{\tuple{d}}/\Delta$ of~$\PG_{\tuple{d}}$.
Its lattice of $1$-parameter subgroups is~$\cocharacters(T) = \cocharacters(T_{\tuple{d}})/\mathbb{Z}\delta$,
where $\delta = \sum_{i\in Q_0}\sum_{r=1}^{d_i}\lambda_{i,r}$,
or more concretely $\delta(z) \!=\! z\cdot \operatorname{id}$.
We may identify~$\cocharacters(T)_\mathbb{R}$ with
the \hbox{orthogonal} complement~$(\mathbb{R}\delta)^\bot$ inside~$\cocharacters(T_{\tuple{d}})_\mathbb{R}$. It~consists of all real linear combinations
\[
\sum_{i\in Q_0}\sum_{r=1}^{d_i} a_{i,r}\lambda_{i,r}
\]
such that $\sum_{i\in Q_0}\sum_{r=1}^{d_i} a_{i,r} = 0$.
We restrict the norm $\|{-}\|$ to this subspace.

Now now recall the terminology for the Harder--Narasimhan stratification.
For a dimension vector $\tuple{e} \!\neq\! 0$ we define $|\tuple{e}| \!=\! \sum_{i \in Q_0} e_i$
and the \emph{slope}~\hbox{$\mu(\tuple{e}) \!=\! \mu_\theta(\tuple{e}) \!\colonequals\! \sfrac{\theta\cdot\tuple{e}}{|\tuple{e}|}$},
i.e., we~choose the total dimension as denominator for the slope.
If~$M$ is a representation of~$Q$, we~write~$\mu(M)\colonequals\mu(\dimvect(M))$ and call it the slope of $M$.

\begin{definition}
Let~$\mu=\mu_\theta$ be as above.
A representation $M$ of $Q$ is called \emph{$\mu$\nobreakdash-semistable} (respectively \emph{$\mu$\nobreakdash-stable})
if any nonzero proper subrepresentation $M'$ satisfies the inequality
\begin{equation*}
\mu(M') \leq \mu(M) \quad(\text{respectively } \mu(M') < \mu(M)).
\end{equation*}
\end{definition}

By clearing the denominator as explained in \cite[Rem.\,2.3]{MR3199484}
this recovers King's definition of (semi)stability, so~the moduli spaces of (semi)stable objects do not depend on the choice of denominator in~$\mu$.
However, the choice of (the denominator in) the slope function
is required to define the Harder--Narasimhan stratification.
We will stick to the standard choice of slope function made in \cite{MR1974891},
and leave it for future work to understand its influence on the results in this paper.

\begin{definition}
Let $M$ be a representation of $Q$.
A \emph{Harder--Narasimhan filtration} of $M$ is a sequence
\begin{equation*}
0 = N^0 \subset N^1 \subset N^2 \subset \ldots \subset N^\ell = M
\end{equation*}
of subrepresentations such that each subquotient $N^m/N^{m-1}$ is $\mu$\nobreakdash-semistable
and such that the chain of inequalities
\begin{equation*}
\mu(N^1/N^0) > \mu(N^2/N^1) > \ldots > \mu(N^\ell/N^{\ell-1})
\end{equation*}
holds.
\end{definition}

\begin{remark}
Any such filtration defines the \emph{Harder--Narasimhan type}
\begin{equation*}
\tuple{d}^* = (\dimvect(N^1/N^0),\ldots,\dimvect(N^{\ell}/N^{\ell-1})).
\end{equation*}
Conversely, for any tuple
$\tuple{d}^* = (\tuple{d}^1,\ldots,\tuple{d}^{\ell})$ summing to $\tuple{d}$ and
such that each $\tuple{d^s}$ is $\mu$\nobreakdash-semistable and
the inequalities $\mu(\tuple{d}^1)>\ldots>\mu(\tuple{d}^{\ell})$ hold,
there exist representations whose Harder--Narasimhan type is~$\tuple{d}^*$.
\end{remark}

Rudakov establishes in \cite[Th.\,3]{MR1480783} the existence
and uniqueness
of the Harder--Nara\-sim\-han filtration,
see also~\cite[Th.\,2.5]{MR1906875} which proves this directly
in the case of quiver representations.

We may stratify~$\repspace{Q,\tuple{d}}$ by Harder--Narasimhan type as follows.
Let~$\tuple{d}^* = (\tuple{d}^1,\ldots,\tuple{d}^\ell)$ be a Harder--Narasimhan type,
i.e.,~a sequence of dimension vectors of strictly decreasing slope
such that for each dimension vector~$\tuple{d}^m$
there exists a~$\mu$\nobreakdash-semistable representation of dimension vector~$\tuple{d}^m$.

For a representation $M \in \repspace{Q,\tuple{d}}$
and~$a\in Q_1$, we~can decompose the matrix
describing the linear transformation~$M_a\colon M_{\source(a)}\to M_{\target(a)}$ into blocks
\begin{equation}
\label{equation:shape-matrix-HN-blocks}
M_a =
\left(
\begin{matrix}
M_a^{1,1} & \hdots & M_a^{1,\ell} \\
\vdots & \ddots & \vdots \\
M_a^{\ell,1} & \hdots & M_a^{\ell,\ell}
\end{matrix}
\right),
\end{equation}
where each block $M_a^{n,m}$ is of size $d^n_{\target(a)} \times d^m_{\source(a)}$,
for~$m,n=1,\ldots,\ell$.
Similarly, each component~$g_i$ of~$g \in \group$
can be decomposed into blocks~$g_i^{n,m}$ of size~$d^n_i \times d^m_i$,
for~$m,n=1,\ldots,\ell$.

Let $L_{\tuple{d}^*}$ and $P_{\tuple{d}^*}$ be subgroups of ${\group}$ defined by
\begin{align*}
L_{\tuple{d}^*} &\colonequals \{ g \in \group \mid g_i^{n,m} = 0 \text{ for all $i \in Q_0$ and all $n \neq m$}\}, \\
P_{\tuple{d}^*} &\colonequals \{ g \in \group \mid g_i^{n,m} = 0 \text{ for all $i \in Q_0$ and all $n > m$}\}.
\end{align*}
The group $L_{\tuple{d}^*}$ is a Levi factor of the parabolic subgroup $P_{\tuple{d}^*}$ of $\group$.
Let~$R_{\tuple{d}^*}$ and~$R_{\tuple{d}^*}^+$ be closed subvarieties of~$R$ defined by
\begin{equation}\label{equation:upper-triangular}
\begin{aligned}
R_{\tuple{d}^*} &\colonequals \{ M \in \repspace{Q,\tuple{d}} \mid M_a^{n,m} = 0 \text{ for all $a \in Q_1$ and all $n \neq m$}\}, \\
R_{\tuple{d}^*}^+ &\colonequals \{ M \in \repspace{Q,\tuple{d}} \mid M_a^{n,m} = 0 \text{ for all $a \in Q_1$ and all $n > m$}\}.
\end{aligned}
\end{equation}
The group~$L_{\tuple{d}^*} \cong \group[\tuple{d}^1] \times \ldots \times \group[\tuple{d}^\ell]$
acts on~$R_{\tuple{d}^*} \cong \repspace{Q,\tuple{d}^1} \times \ldots \times \repspace{Q,\tuple{d}^\ell}$,
and~$P_{\tuple{d}^*}$ acts on~$R_{\tuple{d}^*}^+$.
The projection~$p_{\tuple{d}^*}\colon R_{\tuple{d}^*}^+ \to R_{\tuple{d}^*}$
which forgets the off-diagonal blocks is equivariant with respect to the projection~$P_{\tuple{d}^*} \to L_{\tuple{d}^*}$.

\begin{definition}
\label{definition:harder-narasimhan-strata}
Let $\tuple{d}^* = (\tuple{d}^1,\ldots,\tuple{d}^\ell)$ be a Harder--Narasimhan type.
The locus
\begin{equation*}
R_{\tuple{d}^*}^{\HN} \colonequals \{ M \in \repspace{Q,\tuple{d}} \mid \text{$M$ has a Harder--Narasimhan filtration of type $\tuple{d}^*$}\}
\end{equation*}
is called the associated \emph{Harder--Narasimhan stratum}.

We also define $Z_{\tuple{d}^*} \colonequals \repspace[\mu\semistable]{Q,\tuple{d}^1} \times \ldots \times \repspace[\mu\semistable]{Q,\tuple{d}^\ell}$,
and $\blade[\tuple{d}^*] \colonequals p_{\tuple{d}^*}^{-1}(Z_{\tuple{d}^*})$.
\end{definition}

We state the fifth-named author's result~\cite[Prop.\,3.4]{MR1974891} on the Harder--Nara\-sim\-han stratification.

\begin{theorem}[Reineke]
\label{theorem:HN-stratification}
The affine space $R$ admits a decomposition
\begin{equation*}
R = \bigsqcup_{\tuple{d}^*} R_{\tuple{d}^*}^{\HN},
\end{equation*}
as a finite disjoint union into finitely many locally closed irreducible $\group$-invariant subsets. Moreover,
\begin{enumerate}
\item The set $Z_{\tuple{d}^*}$ is an $L_{\tuple{d}^*}$-invariant open subset of $R_{\tuple{d}^*}$;
\item The set $\blade[\tuple{d}^*]$ is a $P_{\tuple{d}^*}$-invariant open subset of $R_{\tuple{d}^*}^+$;
\item Each stratum $R_{\tuple{d}^*}^{\HN}$ satisfies $R_{\tuple{d}^*}^{\HN} = {\group}\cdot \blade[\tuple{d}^*]$; and
\item \label{item:action-map} The action map~$\sigma\colon\group \times \blade[\tuple{d}^*] \to R_{\tuple{d}^*}^{\HN}$
induces an isomorphism with the associated fiber bundle,
i.e.,~$R_{\tuple{d}^*}^{\HN} \cong \group \times^{P_{\tuple{d}^*}} \blade[\tuple{d}^*]$.
\end{enumerate}
\end{theorem}

We are going to describe how to identify these two stratifications.
Let $\tuple{d}^* = (\tuple{d}^1,\ldots,\tuple{d}^\ell)$ be a Harder--Narasimhan type
and assume that $\tuple{d}^* \neq (\tuple{d})$.
Let $C$ be the minimal positive integer such that
\begin{equation}
\label{equation:ks}
k_m\colonequals C\mu(\tuple{d}^m) \in \mathbb{Z}
\end{equation}
for all $m=1,\ldots,\ell$. In~\cref{table:harder-narasimhan-strata} we have listed the Harder--Narasimhan types
for an interesting Kronecker moduli space,
together with the values of~$\mu(\tuple{d}^m)$, $C$, and~$k_m$
for all types, and all~$m=1,\ldots,\ell$.

We define a $1$-parameter subgroup $\lambda = \lambda_{\tuple{d}^*} = (\lambda_i)_{i \in Q_0} \in \cocharacters(T_{\tuple{d}})$ by
\begin{equation}
\label{equation:HN-type-1-PS-correspondence}
\lambda_i(z) = \operatorname{diag}\big( \underbrace{z^{k_1},\ldots,z^{k_1}}_{d_i^1 \text{ times}}; \underbrace{z^{k_2},\ldots,z^{k_2}}_{d_i^2 \text{ times}};\ldots; \underbrace{z^{k_\ell},\ldots,z^{k_\ell}}_{d_i^\ell \text{ times}} \big).
\end{equation}
Note that $\lambda$ is primitive by the minimality of $C$,
and that
\begin{equation*}
(\delta,\lambda)
=\sum_{i \in Q_0} \sum_{m=1}^\ell k_md_i^m
=\sum_{m=1}^\ell k_m|\tuple{d}^m|
=C\sum_{m=1}^\ell \theta(\tuple{d}^m)
=C\theta(\tuple{d})
= 0.
\end{equation*}
This shows that $\lambda$ lies in $(\mathbb{R}\delta)^\bot$
which we have identified with $\cocharacters(T)_\mathbb{R}$.
We thus can, and will, interpret~$\lambda$ also as a $1$-parameter subgroup of~$\PG_{\tuple{d}}$.

We can now compare the two stratifications obtained in \cref{theorem:hesselink,theorem:HN-stratification}.
The comparison we state next is proved independently in~\cite[Th.\,3.8]{MR3871820} and~\cite[Th.\,6.3]{MR3199484},
where we restrict ourselves to the case where the denominator of the slope function
is the total dimension (and we are not using relations for the quiver).
For a representation~$M$ which has a Harder--Narasimhan filtration of type~$\tuple{d}^*$, we~get $[\lambda] \cap \Lambda(M) \neq \emptyset$.
This shows that every primitive $1$-parameter subgroup~$\lambda \in \cocharacters(T)$
which occurs in the Hesselink stratification for $\repspace{Q,\tuple{d}}$
is of the form~$\lambda_{\tuple{d}^*}$
for a unique Harder--Narasimhan type $\tuple{d}^*$
after conjugation with a suitable Weyl group element.

\begin{theorem}[Hesselink--Harder--Narasimhan correspondence]
\label{theorem:HN=hesselink}
For the Hesselink and Harder--Nara\-sim\-han stratifications
from \cref{theorem:hesselink,theorem:HN-stratification}
\begin{equation*}
R \setminus R^{\theta\semistable} = \bigsqcup_{\tuple{d}^* \neq (\tuple{d})} R_{\tuple{d}^*}^{\HN} = \bigsqcup_{[\lambda]} S_{[\lambda]}
\end{equation*}
the following hold:
\begin{enumerate}
\item For every Harder--Narasimhan type $\tuple{d}^* \neq (\tuple{d})$,
the $1$-parameter subgroup~$\lambda_{\tuple{d}^*}$
satisfies~$\langle \chi_\theta,\lambda_{\tuple{d}^*} \rangle < 0$
and~$R_{\tuple{d}^*}^{\HN} = S_{[\lambda_{\tuple{d}^*}]}$.
\item For ever ${\group}$-conjugacy class $[\lambda]$ in the Hesselink stratification for which~\hbox{$S_{[\lambda]} \mkern-4mu\neq\mkern-4mu \emptyset$},
there exists a unique Harder--Narasimhan type $\tuple{d}^* \neq (\tuple{d})$
such that $[\lambda] = [\lambda_{\tuple{d}^*}]$.
\end{enumerate}
\end{theorem}
The same correspondence also gives rise to several other identifications,
that we state and introduce notation for in the following remark.
\begin{remark}
\label{remark:technicalities-and-abbreviations-HN-Hesselink-equivalence}
Let $\tuple{d}^* \neq (\tuple{d})$ be a Harder--Narasimhan type
and let $\lambda = \lambda_{\tuple{d}^*}$ be the $1$-parameter subgroup of $T\subset\PG_{\tuple{d}}$
which corresponds to it under \cref{theorem:HN=hesselink}.
Then we set
\begin{align*}
L &\colonequals L_{\tuple{d}^*}/\Delta = L_\lambda,&
P &\colonequals P_{\tuple{d}^*}/\Delta = P_\lambda,&
R^0 &\colonequals R_{\tuple{d}^*} = R_\lambda, &
R^+ &\colonequals R_{\tuple{d}^*}^+ = R_\lambda^+,\\
Z &\colonequals Z_{\tuple{d}^*} = Z_\lambda,&
\blade &\colonequals \blade[\tuple{d}^*] = \blade[\lambda], &
S &\colonequals R_{\tuple{d}^*}^{\HN} = S_{[\lambda]}.
\end{align*}
Then we have the following situation:
\begin{equation*}
\begin{tikzcd}
L & P \\[-1.5em]
\curvearrowright & \curvearrowright \\[-1.5em]
R^0 & R^+ \arrow{l}{p} \\[-1.5em]
\text{\rotatebox[origin=c]{90}{$\subseteq$}} & \text{\rotatebox[origin=c]{90}{$\subseteq$}} \\[-1.5em]
Z & \Sigma \arrow{l}{} \arrow{r}{} \arrow{d}{} & S \arrow{d}{} \\
& \{eP\} \arrow{r}{} & {\group}/P.
\end{tikzcd}
\end{equation*}
\end{remark}

\subsection{Width of the windows for quiver moduli}
\label{subsection:width}
Let~$\tuple{d}^*=(\tuple{d}^1,\ldots,\tuple{d}^\ell)$ be a Harder--Narasimhan type,
and let $\lambda \colonequals \lambda_{\tuple{d}^*} \in \characters_*({\group})$ be
the corresponding $1$\nobreakdash-parameter subgroup given by~\eqref{equation:HN-type-1-PS-correspondence}.
Let $S$ be the Hesselink stratum associated to $\lambda\colonequals\lambda_{\tuple{d}^*}$,
or equivalently the Harder--Narasimhan stratum associated to $\tuple{d}^*$,
let $\blade\colonequals\blade[\lambda]$ be the blade of $\lambda$ (and of $\tuple{d}^*$),
and let~$\limitset[]\colonequals\limitset[\lambda]$ be the limit set of $\lambda$ (and of~$\tuple{d}^*$).

We wish to compute the weight of $\det(\normal_{S/R}^\vee|_Z)$
for the action of $\lambda$,
which we denote by $\eta_\lambda$,
which is one of the ingredients in the statement of Teleman quantization as in \cref{theorem:teleman-quantization}. By~using the equivariant adjunction formula,
which we state for a general case below, we~will be able to split the computation in two parts and conduct each of them separately.
\begin{lemma}
\label{lemma:determinant-normal-bundle-sequence}
Let $R$ be a smooth variety,
and let $S$ be a smooth, locally closed subvariety of $R$.
Let a linearly reductive algebraic group $G$ act on $R$,
and assume that~$S$ is $G$-stable. In~this case, the standard adjunction isomorphism
\begin{equation}
\label{equation:equivariant-adjunction}
\det(\normal_{S/R})^\vee \cong \omega_R|_S \otimes \omega_S^\vee
\end{equation}
is $G$-equivariant, i.e., it holds in $\Pic^G(S)$.
\end{lemma}

\begin{proof}
By assumption $S$ is locally closed,
hence closed in an open $U$ of $R$.
Since both $U$ and $S$ are smooth,
the standard adjunction formula applies. This isomorphism is $G$-equivariant because it is induced from the short exact sequence
\begin{equation*}
0 \to \mathcal{I}_{S}/\mathcal{I}^2_{S} \to \Omega^{1}_{U/\field}|_{S} \to \Omega^{1}_{S/\field} \to 0,
\end{equation*}
which is $G$-equivariant as $S$ is $G$-stable.
\end{proof}

The space $R = \repspace{Q,\tuple{d}}$ is defined to be the product
of affine spaces~$\Mat_{d_{\target(a)} \times d_{\source(a)}}$,
indexed by the arrows~$a\in Q_1$, see \eqref{equation:representation-variety}. The coordinate ring of $R$ is thus
\begin{equation*}
\mathcal{O}(R) = \field[x_{p,q}^{(a)}]_{a \in Q_1, ~p=1,\ldots,d_{\target(a)}, ~q=1,\ldots,d_{\source(a)}},
\end{equation*}
where $x_{p,q}^{(a)}$ is the regular function which selects from a tuple $M = (M_a)_{a \in Q_1}$ the $(p,q)$-th entry of the matrix $M_a$.
The sheaf of differentials $\Omega^{1}_R$ is therefore the sheaf associated to the free $\mathcal{O}(R)$-module
with basis $\smash{\{\diff x_{p,q}^{(a)}\}}$ where~$a \in Q_1,p=1,\ldots,d_{\target(a)}$ and~$q=1,\ldots,d_{\source(a)}$.

The action of $\group$ on $R\colonequals\repspace{Q,\tuple{d}}$ induces a left action on the coordinate ring, by~precomposition with the inverse, so~for $f \in \mathcal{O}(R)$ and $g = (g_i)_{i \in Q_0} \in \group$,
the regular function $g\cdot f$ is defined by
\begin{equation*}
(g\cdot f)(M) = f(g^{-1}\cdot M) = f((g_{\target(a)}^{-1}M_ag_{\source(a)})_{a \in Q_1})
\end{equation*}
for all $M = (M_a)_{a \in Q_1} \in R$.
The induced left action of $\group$ on the sheaf of differentials~$\Omega^1_R$
is given by~$g \cdot \mathrm{d}f \colonequals \mathrm{d}(g \cdot f)$, for all $f \in \mathcal{O}(R)$.
This yields left actions on all exterior products and all $\group$-stable subsheaves of $\Omega^1_R$, as~well as on quotients of $\Omega^1_R$ by such subsheaves.

First we deal with the first tensor factor in \eqref{equation:equivariant-adjunction}
restricted to~$Z$,
i.e., we~consider the limit set~$Z$ inside stratum~$S$
which itself lives inside the representation variety~$R$.
\begin{lemma}
\label{lemma:omega-RZ-weight}
The $\lambda$-weight of the canonical bundle~$\omega_{R}|_Z$ on~$R$ restricted is
\begin{equation}
%\label{equation:omega-RZ-weight}
\weight_{\lambda}(\omega_R|_Z) = \sum_{1\leq m<n\leq\ell}(k_n-k_m)\left( \langle \tuple{d}^m,\tuple{d}^n\rangle - \langle \tuple{d}^n,\tuple{d}^m\rangle \right).
\end{equation}
\end{lemma}

\begin{proof}
Recall that $\lambda = (\lambda_i)_{i \in Q_0}$ consists of diagonal $1$-parameter subgroups $\lambda_i$,
see~\eqref{equation:HN-type-1-PS-correspondence}.
The action of~$\lambda$ on~$x_{p,q}^{(a)}$ is given by
\begin{equation*}
(\lambda(z) \cdot x_{p,q}^{(a)})(M) = x_{p,q}^{(a)}((\lambda_{\target(b)}(z)^{-1}M_b\lambda_{\source(b)}(z))_{b \in Q_1})
= \lambda_{\target(a)}(z)_{p,p}^{-1}\lambda_{\source(a)}(z)_{q,q}\cdot x_{p,q}^{(a)}(M)
\end{equation*}
for all $M = (M_b)_{b \in Q_1}$ and all $z \in \Gm$.
This shows that
\begin{equation*}
\lambda(z)\cdot \diff x_{p,q}^{(a)} = \lambda_{\target(a)}(z)_{p,p}^{-1}\lambda_{\source(a)}(z)_{q,q} \diff x_{p,q}^{(a)}.
\end{equation*}
Taking the exterior product over all the generators $\diff x_{p,q}^{(a)}$, we~obtain
\begin{equation*}
\bigwedge_{a \in Q_1}\bigwedge_{\substack{1 \leq p \leq d_{\target(a)} \\ 1 \leq q \leq d_{\source(a)}}} \!\! \lambda(z)\cdot \diff x_{p,q}^{(a)}
= \Biggl( \prod_{a \in Q_1} \prod_{\substack{1 \leq p \leq d_{\target(a)} \\ 1 \leq q \leq d_{\source(a)}}} \!\! \lambda_{\target(a)}(z)_{p,p}^{-1}\lambda_{\source(a)}(z)_{q,q} \Biggr) \bigwedge_{a \in Q_1} \bigwedge_{\substack{1 \leq p \leq d_{\target(a)} \\ 1 \leq q \leq d_{\source(a)}}} \!\!\diff x_{p,q}^{(a)}.
\end{equation*}
This determines the character of the $\Gm$-linearization via $\lambda$ of $\omega_R$.
Its weight is necessarily $\weight_\lambda(\omega_R|_Z)$. So, we~obtain
\begin{align*}
\weight_\lambda(\omega_R|_Z)
&= \sum_{a \in Q_1} \sum_{1 \leq m,n \leq \ell} (k_m - k_n) d_{\source(a)}^m d_{\target(a)}^n \\
&= \sum_{a \in Q_1} \sum_{1 \leq m < n \leq \ell}(k_m - k_n)(d_{\source(a)}^m d_{\target(a)}^n - d_{\source(a)}^n d_{\target(a)}^m) \\
&= \sum_{1 \leq m < n \leq \ell}(k_m - k_n)(\langle \tuple{d}^n,\tuple{d}^m \rangle - \langle \tuple{d}^m,\tuple{d}^n \rangle) \\
&= \sum_{1 \leq m < n \leq \ell}(k_n - k_m)(\langle \tuple{d}^m,\tuple{d}^n \rangle - \langle \tuple{d}^n,\tuple{d}^m \rangle).\qedhere
\end{align*}
\end{proof}

Next we deal with the second tensor factor in \eqref{equation:equivariant-adjunction}
restricted to~$Z$,
i.e., we~consider the limit set~$Z$ inside the stratum~$S$,
without reference to~$R$.
\begin{lemma}
\label{lemma:omega-SZ-weight}
The~$\lambda$-weight of~$\omega_S|_Z$ is
\begin{equation}
%\label{equation:omega-SZ-weight}
\weight_{\lambda}(\omega_S|_Z) = \sum_{1\leq m<n\leq\ell}(k_m-k_n)\langle \tuple{d}^n,\tuple{d}^m\rangle.
\end{equation}
\end{lemma}
To give a proof, we~first explain how to split the computation in several steps,
eventually leading to the identity in~\eqref{equation:split-computation},
and then perform each step separately.

The morphism~$\sigma\colon \group \times \blade \to S$
in \cref{theorem:HN-stratification}(\ref{item:action-map})
is a principal fiber bundle (for the étale topology, by~the standing assumption on~$\field$ being of characteristic zero)
with fiber $P_{\tuple{d}^*}$.
Therefore, $\sigma$ is smooth and the relative tangent bundle sequence
\begin{equation}
\label{equation:relative-tangent-bundle-sequence}
0
\to\tangent_{\group \times \blade/S}
\to\tangent_{\group \times \blade}
\to\sigma^*\tangent_{S}
\to0.
\end{equation}
is exact.
The fiber of $\sigma$ in a point $\sigma(g,M) = g\cdot M$ is the $P_{\tuple{d}^*}$-orbit of $(g,M)$.
The latter is isomorphic to~$P_{\tuple{d}^*}$ via the action of $P_{\tuple{d}^*}$ on the point $(g,M)$, which is given~by
\begin{equation*}
p \cdot (g,M) = (gp^{-1}, p\cdot M).
\end{equation*}
Denote~$\mathfrak{g}_{\tuple{d}}\colonequals\operatorname{Lie}\group$
and~$\mathfrak{p}_{\tuple{d}^*}\colonequals\operatorname{Lie}P_{\tuple{d}^*}$
the associated Lie algebras.
With this notation,
the sequence \eqref{equation:relative-tangent-bundle-sequence} can be written as
\begin{equation}
\label{equation:relative-tangent-bundle-sequence-rewritten}
0
\to\mathfrak{p}_{\tuple{d}^*} \otimes \mathcal{O}_{\group \times \blade}
\to\tangent_{\group \times \blade}
\to\sigma^*\tangent_{S}
\to0,
\end{equation}
where, over a point $(g,M) \in \group \times \blade$,
the map $\mathfrak{p}_{\tuple{d}^*} \to T_{\group \times \blade,(g,M)} = g\cdot \mathfrak{g}_{\tuple{d}} \oplus T_{\blade,M}$
is the derivative of the map $P_{\tuple{d}^*} \to \group \times \blade$
defined by $p \mto p\cdot (g,M) = (gp^{-1}, p\cdot M)$.

The sequence \eqref{equation:relative-tangent-bundle-sequence}
is $P_{\tuple{d}^*}$-equivariant, which implies that~\eqref{equation:relative-tangent-bundle-sequence-rewritten} is also $P_{\tuple{d}^*}$\nobreakdash-equi\-variant;
note that here we consider the adjoint action of~$P_{\tuple{d}^*}$ on~$\mathfrak{p}_{\tuple{d}^*}$, which is induced by the conjugation action of $P_{\tuple{d}^*}$ on $\group$, the left multiplication on $\Sigma$ and the action induced by $\sigma$ on~$\tangent_{S}$.

Now we restrict the action of $P_{\tuple{d}^*}$ via the $1$-parameter subgroup $\lambda\colon \Gm\!\to\! P_{\tuple{d}^*} \!\subseteq\! \group$.
As $Z$ is the locus of fixed points of the $\lambda$-action,
$\lambda$ acts on every fiber of $T_S|_Z$,
and as $Z$ is connected, the action is the same on every fiber.
Let $M \in Z$ and consider the sequence \eqref{equation:relative-tangent-bundle-sequence-rewritten}
in the fiber of the point~$(e,M) \in \group \times \blade$. It~is
\begin{equation}
\label{equation:relative-tangent-bundle-sequence-fiber}
0
\to\mathfrak{p}_{\tuple{d}^*}
\to\mathfrak{g}_{\tuple{d}} \oplus \tangent_{\blade,M}
\to\tangent_{S,M}
\to0.
\end{equation}
The point $(e,M)$ is a fixed point for the $\lambda$-action on $\group \times \blade$,
which implies that $\eqref{equation:relative-tangent-bundle-sequence-fiber}$ is a short exact sequence of representations of $\Gm$.

As the blade $\blade$ is open in the affine space $R_{\tuple{d}^*}^+$, see \eqref{equation:upper-triangular},
the tangent space to $\blade$ at every point is identified
with $R_{\tuple{d}^*}^+$ (considered as a vector space).
The exact sequence \eqref{equation:relative-tangent-bundle-sequence-fiber} thus becomes
\begin{equation}
\label{equation:relative-tangent-bundle-sequence-fiber-rewritten}
0
\to\mathfrak{p}_{\tuple{d}^*}
\to\mathfrak{g}_{\tuple{d}} \oplus R_{\tuple{d}^*}^+
\to\tangent_{S,M}
\to0.
\end{equation}
Although it is not relevant for the weight computation, we~can describe the map $\mathfrak{p}_{\tuple{d}^*} \to \mathfrak{g}_{\tuple{d}} \oplus R_{\tuple{d}^*}^+$ explicitly. It~sends $x \in \mathfrak{p}_{\tuple{d}^*}$ to $(-x, [x,M])$, where
\begin{equation*}
[x,M] = (x_{\target(a)}M_a - M_ax_{\source(a)})_{a \in Q_1}.
\end{equation*}
We use the sequence \eqref{equation:relative-tangent-bundle-sequence-fiber-rewritten}
to compute the weight of the restriction to $Z$ of the anticanonical bundle of $S$. It~is
\begin{equation}
\label{equation:split-computation}
\weight_{\lambda}(\omega_S^\vee|_Z) = \weight_{\lambda}(\det(\mathfrak{g}_{\tuple{d}})) + \weight_{\lambda}(\det(R_{\tuple{d}^*}^+)) - \weight_{\lambda}(\det(\mathfrak{p}_{\tuple{d}^*})).
\end{equation}

Now to the individual weights in the right-hand side of~\eqref{equation:split-computation}.
The $1$-parameter subgroup~$\lambda$ acts on $\mathfrak{g}_{\tuple{d}}$ by conjugation
and $\mathfrak{p}_{\tuple{d}^*}$ is a submodule.
On $R_{\tuple{d}^*}^+$,
it acts by restriction of the $\group$-action via~$\lambda\colon \Gm \to \group$.
As~$\lambda$ consists of diagonal matrices,
the matrix entries of elements of~$\mathfrak{g}_{\tuple{d}}$
and of~$R_{\tuple{d}^*}^+$ are weight spaces.
Using the decomposition of the summands of $R_{\tuple{d}^*}^+$
into blocks as in \eqref{equation:shape-matrix-HN-blocks}
and similarly for the summands of $\mathfrak{g}_{\tuple{d}}$ and $\mathfrak{p}_{\tuple{d}^*}$, we~obtain the following.
Here, $\field(r)$ is the one-dimensional~$\Gm$\dash representation
whose weight is $r$.
\begin{lemma}
\label{lemma:weights-fibers}
Regarded as $\Gm$-representations via $\lambda$, we~have the following isomorphisms:
\begin{align*}
\mathfrak{g}_{\tuple{d}} &\cong \bigoplus_{i \in Q_0} \bigoplus_{1 \leq m,n \leq \ell} \field(k_m-k_n)^{d_i^md_i^n} \\
\mathfrak{p}_{\tuple{d}^*} &\cong \bigoplus_{i \in Q_0} \bigoplus_{1 \leq m<n \leq \ell} \field(k_m-k_n)^{d_i^md_i^n} \\
R_{\tuple{d}^*}^+ &\cong \bigoplus_{a \in Q_1} \bigoplus_{1 \leq m < n \leq \ell} \field(k_m-k_n)^{d_{\target(a)}^md_{\source(a)}^n}
\end{align*}
\end{lemma}
The above lemma implies at once the following.
\begin{lemma}
\label{lemma:weights-determinants}
The $\lambda$-weights of the determinants of $\mathfrak{g}_{\tuple{d}}$, $\mathfrak{p}_{\tuple{d}^*}$, and $R_{\tuple{d}^*}^+$ are
\begin{align*}
\weight_\lambda(\det(\mathfrak{g}_{\tuple{d}})) &= 0 \\
\weight_\lambda(\det(\mathfrak{p}_{\tuple{d}^*})) &= \sum_{1 \leq m < n \leq \ell} (k_m - k_n) \Bigl(\sum_{i \in Q_0} d_i^n d_i^m\Bigr) \\
\weight_\lambda(\det(R_{\tuple{d}^*}^+)) &= \sum_{1 \leq m < n \leq \ell} (k_m - k_n) \Bigl(\sum_{a\in Q_1} d_{\target(a)}^m d_{\source(a)}^n\Bigr).
\end{align*}
\end{lemma}

We can now give the proof of \cref{lemma:omega-SZ-weight}.
\begin{proof}[Proof of \cref{lemma:omega-SZ-weight}]
Using~\eqref{equation:split-computation}, we~conclude that the $\lambda$-weight of $\omega_S^\vee|_Z$ is
\begin{align*}
\weight_{\lambda}(\omega_S^\vee|_Z)
&= \sum_{1 \leq m < n \leq \ell} (k_m - k_n) \Bigl(\sum_{a \in Q_1} d_{\source(a)}^n d_{\target(a)}^m - \sum_{i \in Q_0} d_i^n d_i^m\Bigr) \\
&= \sum_{1\leq m < n \leq \ell} (k_n - k_m) \langle \tuple{d}^n,\tuple{d}^m \rangle.\qedhere
\end{align*}
\end{proof}
Thus from \cref{lemma:determinant-normal-bundle-sequence}
and the computations in \cref{lemma:omega-RZ-weight,lemma:omega-SZ-weight}
we obtain the first ingredient to apply Teleman quantization.
\begin{corollary}
\label{corollary:determinant-weight}
The $\lambda$-weight,
where~$\lambda=\lambda_{\tuple{d}^*}$,
of~$\det\bigl(\normal_{S/R}^\vee|_Z\bigr)$ is
\begin{equation}
%\label{equation:weight-eta-notation}
\eta_\lambda\colonequals\weight_{\lambda}(\det\normal_{S/R}^\vee|_Z) = \sum_{1 \leq m < n \leq \ell}(k_n - k_m)\langle \tuple{d}^m,\tuple{d}^n \rangle.
\end{equation}
\end{corollary}
Throughout we will denote this weight by $\eta_\lambda$. In~\cref{table:harder-narasimhan-strata} we have computed~$\eta_\lambda$
for all Harder--Narasimhan strata,
in our running example of an interesting Kronecker moduli space.

\subsection{Weights of the endomorphism bundle of the universal bundle}
\label{subsection:weights-endomorphism-bundle}
Recall from \cref{section:moduli} that,
if $\tuple{d}$ is $\theta$-coprime,
the moduli space~$\modulispace[\theta\stable]{Q,\tuple{d}}$
is fine, i.e.,
comes equipped with a universal bundle~$\mathcal{U}=\mathcal{U}(\tuple{a})$.
As explained in \cref{remark:uniqueness-up-to-twist-universal-bundle},
this bundle is unique up to the choice of a normalization,
which is given by a tuple~$\tuple{a}\in\mathbb{Z}^{Q_0}$ such that~$\tuple{a}\cdot\tuple{d}=1$.

The bundles~$\mathcal{U}_i(\tuple{a})$ are the descent of
(the restriction to the stable locus of)
the~$\group$\nobreakdash-equivariant bundles~$U_i(\tuple{a})=U_i\otimes L(\tuple{a})$
on~$R=\repspace{Q,\tuple{d}}$.
We will now consider these equivariant bundles on~$R$ itself,
thus in what follows, we~do not have to put any conditions to ensure their existence,
nor do we have to choose a normalization with certain properties.

For a Harder--Narasimhan type $\tuple{d}^*$, we~will compute the weights of the action of $\lambda_{\tuple{d}^*}$
on the universal bundles $U_i(\tuple{a})$ on $R$.

Recall that the element $g \in {\group}$ acts on $u_i \in U_i(\tuple{a})$ as
\begin{equation*}
g \cdot u_i \colonequals \Bigl(\prod_{j \in Q_0} \det(g_j)^{-a_j} \Bigr) g_i u_i.
\end{equation*}
We have defined the integers~$k_m$ for~$m=1,\ldots,\ell$ in~\eqref{equation:ks}
as the smallest integer multiples of~$\mu(\tuple{d}^m)$,
and used it to define the~$1$-parameter subgroup $\lambda_{\tuple{d}^*}$ in~\eqref{equation:HN-type-1-PS-correspondence}.
From the block decomposition in~\eqref{equation:HN-type-1-PS-correspondence}
we obtain that~$z\in\lambda=\lambda_{\tuple{d}^*}$ acts by
\begin{equation}
\label{equation:1-PS-acting-on-Ui(a)}
z \cdot u_i
= \Bigl(\prod_{j\in Q_0}\det(\lambda_j(z))^{-a_j}\Bigr)\lambda_i(z)u_i
= \Bigl(\prod_{j\in Q_0} z^{-a_j\sum_{n=1}^\ell d^n_j k_n}\Bigr) \lambda_i(z) u_i.
\end{equation}

\begin{lemma}
The weights of the action of~$\lambda=\lambda_{\tuple{d}^*}$
on $U_i(\tuple{a})$ are
\begin{equation}
\label{equation:weights-Ui(a)}
\biggl\{ k_m - \sum_{j \in Q_0} \sum_{n=1}^\ell a_j d^n_j k_n \Bigm| 1\leq m \leq \ell \biggr\},
\end{equation}
where the weight indexed by~$m$ appears with multiplicity~$d_i^m$.
\end{lemma}

\begin{proof}
It suffices to observe in~\eqref{equation:1-PS-acting-on-Ui(a)}
that the weights of~$\lambda_i(z)$ acting on~$u_i$ are~$k_1,\ldots,k_\ell$ and the dimension of the weight space of weight $k_m$ is $d_i^m$.
\end{proof}

The choice of normalization disappears when we consider the summands of the endomorphism bundle, as~in the statement of the following proposition.

\begin{proposition}
\label{proposition:weights-endomorphism-summands}
Let~$Q$ be a quiver,
$\tuple{d}$ a dimension vector,
and~$\theta$ a stability parameter.
Let~$\tuple{d}^*$ be a Harder--Narasimhan type.

The weights of the action of~$\lambda_{\tuple{d}^*}$
on the $\group$-equivariant vector bundles~$U_i^\vee\otimes U_j$
are given by
\begin{equation*}
\{ k_m - k_n\mid 1 \leq m,n \leq \ell \},
\end{equation*}
where~$k_m$ is defined as in~\eqref{equation:ks},
and the weight~$k_m-k_n$ appears with multiplicity~$d_i^md_j^n$.
\end{proposition}

\begin{proof}
Observe that~$U_i^\vee\otimes U_j=U_i^\vee\otimes L(-\tuple{a})\otimes U_j\otimes L(\tuple{a})$
so that it suffices to cancel the contribution of~$L(\tuple{a})$ in~\eqref{equation:weights-Ui(a)}
and sum them together with opposite signs.
\end{proof}

\section{Main theorem and applications}
\label{section:results}
After the work in \cref{section:teleman-for-quiver-moduli}
we can prove the main results in this paper.
For this we introduce in \cref{subsection:strong-ample-stability}
a strengthening of the notion of
an amply stable dimension vector (with respect to some~$\theta$),
which allows us to prove \cref{proposition:inequality},
the main inequality for the proof of \cref{theorem:cohomology-vanishing}.

\subsection{Strong ample stability}
\label{subsection:strong-ample-stability}
We introduce a strengthening of the notion of being amply~$\theta$-stable.
The former was introduced in \cite[Def.\,4.1]{MR3683503}
and recalled in \cref{definition:ample-stability},
the strengthening was introduced without giving it a name in \cite[Prop.\,5.1]{MR3683503}, as~a sufficient numerical condition for ample stability.
\begin{definition}
\label{definition:strong-ample-stability}
The dimension vector $\tuple{d}$ is said to be \emph{strongly amply~$\theta$-stable}
if for any dimension subvector $\tuple{e}$ for which $\mu(\tuple{e}) > \mu(\tuple{d} - \tuple{e})$,
(or equivalently by clearing denominators,~$\theta(\tuple{e})>0$)
the inequality $\langle \tuple{e}, \tuple{d} - \tuple{e} \rangle \leq -2$ holds.
\end{definition}
Clearing the denominator shows that strong ample stability only depends on~$\theta$,
not on the choice of denominator for~$\mu$.
The condition above has been shown~\cite[Prop.\,5.1]{MR3683503} to imply ample stability.
As we will see in \cref{example:strict-ample-stability}, the converse is not true.
Strong ample stability has also been studied by Martinez Acosta in \cite{karen},
where it goes by the name of ``numerically amply $\theta$-stable'',

\begin{remark}
We can interpret this definition as saying that
the Harder--Nara\-sim\-han stratum
associated to the 2-step Harder--Nara\-sim\-han type~$\tuple{d}^*=(\tuple{e},\tuple{d}-\tuple{e})$,
has codimension $-\langle \tuple{e}, \tuple{d} - \tuple{e} \rangle\geq 2$,
provided it is non-empty.
\end{remark}

We will recall some positive results from \cite{MR3683503,karen},
and subsequently prove a general criterion,
indicating that strong ample stability is a natural condition. It~can also be algorithmically checked,
and it is one of the features of \textsc{QuiverTools} \cite{quivertools},
see also~\cite{2506.19432}.

A negative result,
where strong ample stability does not hold,
whilst ample stability does,
is given in \cref{example:strict-ample-stability}. It~was the smallest
(and somewhat artificial) example
we could find,
suggesting strong ample stability is indeed satisfied in many examples of interest
where ample stability holds.

First some special cases of often studied quivers.
\begin{lemma}
\label{lemma:strong-ample-stability-examples}
Strong ample stability holds for the following cases:
\begin{itemize}
\item the $m$-Kronecker quiver with~$m\geq 3$,
and~$\tuple{d}$ a stable dimension vector
(in which case one can assume that~$\theta=\theta_\can$),
except for~$\tuple{d}=(2,2)$.
\item the generalized Jordan quiver (or loop quiver)
with~$m\geq 2$ loops
and~$\tuple{d}=(d)$ a~dimension vector with~$d\geq 2$,
except for~$(m,d)=(2,2)$.
\item the $m$-subspace quiver (so~$\#Q_0=m+1$) with~$m\geq 3$,
$\tuple{d}$ of the form~$(1,\ldots,1,d)$ with~$d\neq m-1$
which is~$\theta_\can$-coprime.
\end{itemize}
\end{lemma}

\begin{proof}
For the Kronecker case it follows from the proof of \cite[Prop.\,6.2]{MR3683503}.
The generalized Jordan quiver is covered on \cite[p.\,460]{MR3683503}.
The subspace quiver case is (a~special case of) \cite[Th.\,3.3.4]{karen}.
\end{proof}

The two exceptional cases excluded in \cref{lemma:strong-ample-stability-examples}
are not amply stable.
We can often vary the stability parameter,
and preserve strong ample stability, as~illustrated by this geometrically meaningful example.
\begin{example}
Let~$Q$ be the~6-subspace quiver,
and let~$\tuple{d}=(1,1,1,1,1,1;2)$. By~\cref{lemma:strong-ample-stability-examples},
this setting is strongly amply stable for the canonical stability parameter~$\theta_\can$.
However, this data gives rise to a singular moduli space \cite{MR3778138,2202.08601}.

Two non-isomorphic small resolutions of singularities have been studied in \cite{MR3778138}, by~considering the stability parameters
\begin{equation*}
\begin{aligned}
\theta^+&=\Bigl( \frac{1}{3}+\epsilon,\frac{1}{3}-\frac{\epsilon}{5},\ldots,\frac{1}{3}-\frac{\epsilon}{5};-1 \Bigr) \\
\theta^-&=\Bigl( \frac{1}{3}-\epsilon,\frac{1}{3}+\frac{\epsilon}{5},\ldots,\frac{1}{3}+\frac{\epsilon}{5};-1 \Bigr)
\end{aligned}
\end{equation*}
for some small~$\epsilon>0$
from Section~3~in op.~cit.
One can verify, e.g., using \cite{quivertools},
that~$\tuple{d}$ is strongly amply stable for both stability parameters.
\end{example}

The following proposition gives a general procedure for
finding strongly amply stable dimension vectors
for the canonical stability condition~$\theta_\can=\{\tuple{d},-\}=\langle\tuple{d},-\rangle-\langle-,\tuple{d}\rangle$.
We will say that a quiver~$Q$ and a dimension vector~$\tuple{d}$
have a \emph{thin bridge}
if there exists a decomposition~$Q_0=I'\sqcup I''$
with a unique arrow~$a\colon i'\to i''$ connecting~$I'$ and~$I''$
(so~$i'\in I'$ and~$i''\in I''$)
so that~$d_{i'}=d_{i''}=1$.
We recall that the symmetric bilinear form~$(-,-)$ is defined to be
the symmetrization of the Euler form, i.e.,~$(\alpha,\beta) = \langle \alpha , \beta\rangle + \langle \beta,\alpha\rangle$.
\begin{proposition}
\label{proposition:strong-ample-stability-fundamental-domain}
Let~$Q$ be an acyclic quiver
and let~$\tuple{d}$ be a sincere dimension vector,
i.e., $d_i\geq 1$ for all~$i\in Q_0$.
Assume that
\begin{enumerate}
\item $\tuple{d}$ is~$\theta_\can$-coprime;
\item $\tuple{d}$ is in the interior of the fundamental domain,
i.e., for all~$i\in Q_0$ we have that~$(\tuple{d},\tuple{i})\leq -1$; and
\item $Q$ has no thin bridge.
\end{enumerate}
Then~$\tuple{d}$ is strongly amply~$\theta_\can$-stable.
\end{proposition}
As recalled in \cref{section:moduli}, by~\cite{MR4352662}
the canonical stability~$\theta_\can$ has the pleasant property
that its associated moduli spaces are Fano varieties,
thus making this stability parameter a natural choice.
We relegate the proof of \cref{proposition:strong-ample-stability-fundamental-domain}
to the appendix because it is rather lengthy
and the methods are not related to the rest of the paper.

An interesting corollary to this general procedure is the following result for $3$-vertex quivers.

\begin{corollary}
\label{corollary:hans}
Let~$Q$ be a connected acyclic $3$-vertex quiver,
and let~$\tuple{d}$ be a sincere dimension vector in the fundamental domain
which is~$\theta_\can$-coprime.
Then it is strongly amply~$\theta_\can$-stable.
\end{corollary}

\begin{proof}
By \cref{proposition:strong-ample-stability-fundamental-domain}
we only need to consider $3$-vertex quivers with thin bridges.
There are precisely~6 possibilities,
which have the same symmetrized Euler form up to conjugation by a permutation matrix, so~it suffices to consider one case.
Consider the $3$-vertex quiver on the vertices~$\{1,2,3\}$,
with a thin bridge between~$1$ and~$2$,
and~$m\geq 1$ arrows from~$2$ to~$3$
(and no arrows from~$1$ to~$3$),
and dimension vector~$(1,1,d)$ for some~$d\geq 1$.
The symmetrized Euler form is given by
\begin{equation*}
C=
\begin{pmatrix}
2 & -1 & 0 \\
-1 & 2 & -m \\
0 & -m & 2
\end{pmatrix}
\end{equation*}
so that~$(\tuple{d},\mathbf{1})=(1,1-m,-m+2d)(1,0,0)^{\mathrm{T}}=1$,
and thus~$\tuple{d}$ is not in the fundamental domain.
\end{proof}

\begin{remark}
The special case of the $3$-vertex quiver of the form
\begin{equation}
%\label{equation:karen-quiver}
\begin{tikzpicture}[baseline = -3pt, node distance = 1.5cm]
\node (1) {};
\node (2) [right of = 1] {};
\node (3) [right of = 2] {};
\draw (1) circle (2pt);
\draw (2) circle (2pt);
\draw (3) circle (2pt);
\draw[->, bend left = 30] (1) edge (2);
\draw[->, bend right = 30] (1) edge (2);
\draw[->, bend left = 30] (3) edge (2);
\draw[->, bend right = 30] (3) edge (2);
\end{tikzpicture}
\end{equation}
is studied in \cite[Th.\,3.3.2]{karen}
more exhaustively.
The dimension vectors~$(1,a+1,a)$ and~$(a,a+1,1)$,
mentioned in op.~cit.,
are not~$\theta_\can$-coprime
(it suffices to consider~$\tuple{e}=(1,1,0)$ \resp $(0,1,1)$)
so \cref{corollary:hans} does not apply to them.
\end{remark}

For more examples one can consult \cite[\S3]{karen}.
The next example shows that strong ample stability
is a strictly stronger notion than ample stability, thus justifying its name.
We will revisit it in \cref{example:strict-ample-stability-rigid}
to explain that \cref{theorem:cohomology-vanishing} still holds,
without using \cref{proposition:inequality}.

\begin{example}
\label{example:strict-ample-stability}
Let $Q$ be the $3$-vertex quiver
\begin{equation}
%\label{equation:3-vertex-quiver-1}
Q\colon
\begin{tikzpicture}[baseline = -20pt, node distance = 1.5cm]
\node (1) {};
\node (2) [right of = 1] {};
\node (3) [below of = 2] {};
\draw (1) circle (2pt) node[above] {1};
\draw (2) circle (2pt) node[above] {2};
\draw (3) circle (2pt) node[below] {3};
\draw[->] (1) edge (3);
\draw[->] (2) edge (3);
\draw[->, bend left = 30] (1) edge (2);
\draw[->, bend left = 15] (1) edge (2);
\draw[->] (1) edge (2);
\draw[->, bend right = 15] (1) edge (2);
\draw[->, bend right = 30] (1) edge (2);
\end{tikzpicture}
\end{equation}
Let~$\tuple{d} = (4,1,4)$,
and let $\theta = \theta_\can = (9,-16,-5)$ be the canonical stability parameter.
The resulting moduli space is a Fano 8-fold,
with even Betti numbers given by~$1,2,3,4,5,4,3,2,1$.

The dimension vector $\tuple{d}$ is $\theta$-coprime,
and the stable locus $\repspace[\theta\stable]{Q,\tuple{d}}$ is non-empty.
The representation variety admits a Harder--Narasimhan stratification with 41~strata:
40~unstable strata, and the dense stratum of~$\theta$\dash stable representations.
These can easily be verified using \cite{quivertools}.

Via the formula of~\cite[Prop.\,3.4]{MR1974891}, we~can verify that the codimension of the unstable locus is~$2$, therefore ample stability holds.
On the other hand, strong ample stability is not satisfied:
the dimension vector $\tuple{e} = (3,1,2)$ gives the inequality~$\sfrac{1}{6} = \mu(\tuple{e}) > \mu(\tuple{d-e}) = \sfrac{-1}{3}$,
but we have $\langle \tuple{e}, \tuple{d} - \tuple{e}\rangle = -1$.
The associated 2-step Harder--Narasimhan type~$(\tuple{e},\tuple{d}-\tuple{e})$
does not appear in the list of Harder--Narasimhan types for this setup, as~can be checked using~\cite{quivertools},
which is why it does not influence the codimension of the unstable locus.

In view of \cref{corollary:hans}, we~also point out that~$\tuple{d}$ is not in the fundamental domain,
because~$(\tuple{d},\mathbf{3})=3$.

We revisit this example in \cref{example:strict-ample-stability-rigid}.
\ifcode
Corresponding Sage code below. Version of quiver.tools: commit a07f4f0

\begin{verbatim}
sage: from quiver import *
sage: Q = ThreeVertexQuiver(5,1,1)
sage: d = vector([4,1,4])
sage: theta = Q.canonical_stability_parameter(d)
sage: theta
(9, -16, -5)
sage: is_coprime_for_stability_parameter(d,theta)
True
sage: Q.has_semistable_representation(d,theta)
True
sage: len(Q.all_harder_narasimhan_types(d,theta))
41
sage: 1 - Q.euler_form(d,d)
8
sage: Q.does_rigidity_inequality_hold(d,theta)
True
sage: ample_stability = True
....: for hntype in Q.all_harder_narasimhan_types(d,theta):
....: codim = 0
....: if hntype != [d]:
....: for i in range(0,len(hntype)-1):
....: for j in range(i+1,len(hntype)):
....: codim = codim - Q.euler_form(hntype[i], hntype[j])
....: if codim < 2:
....: ample_stability = False
....:
sage: ample_stability
True
\end{verbatim}
However,
\begin{verbatim}
sage: e = vector([3,1,2])
sage: slope(e,theta)
1/6
sage: slope(d-e,theta)
-1/3
sage: Q.euler_form(e,d-e)
-1
\end{verbatim}
\fi
\end{example}

\subsection{Cohomology vanishing and Schofield's conjecture}
\label{subsection:cohomology-vanishing}
We recall the notation we will use.
For any Harder--Narasimhan type $\tuple{d}^* = (\tuple{d}^1,\dots,\tuple{d}^\ell)$,
let $\lambda \colonequals \lambda_{\tuple{d}^*}$ be the corresponding $1$-parameter subgroup
given in~\eqref{equation:HN-type-1-PS-correspondence};
and let\vspace*{-3pt}
\begin{equation*}
k_m \colonequals C\mu(\tuple{d}^m)=C\theta(\tuple{d}^m)/\lvert\tuple{d}^m\rvert
\end{equation*}
be the smallest integer multiple of the slopes
as in~\eqref{equation:ks},
and let\vspace*{-3pt}
\begin{equation*}
\eta_\lambda\colonequals\weight_\lambda(\det\normal_{S/R}^\vee|_Z)
\end{equation*}
as in \cref{corollary:determinant-weight}.
Recall also the result of \cref{proposition:weights-endomorphism-summands}
describing the weights of~$U_i^\vee\otimes U_j$.

\begin{proposition}
\label{proposition:inequality}
Let~$Q$ be a quiver,
$\tuple{d}$ a dimension vector,
and~$\theta$ a stability parameter,
such that~$\tuple{d}$ is strongly amply~$\theta$-stable.
Let~$\lambda=\lambda_{\tuple{d}^*}$ be the $1$-parameter subgroup
corresponding to a Harder--Narasimhan type~$\tuple{d}^*$.
For all $1 \leq m,n\leq\ell$, the inequality\vspace*{-3pt}
\begin{equation}
\label{equation:teleman-inequality}
k_m-k_n<\eta_\lambda
\end{equation}
holds.
\end{proposition}

Observe that\vspace*{-5pt}
\begin{equation*}
k_1-k_\ell=\max_{1\leq m<n\leq\ell}k_m-k_n,
\end{equation*}
because the entries of the tuple~$(k_m)_{m=1,\ldots,\ell}$ are strictly decreasing by definition.
Moreover,
rewriting~$k_m$ in terms of the slope
(using its definition in~\eqref{equation:ks})
and spelling out~$\eta_\lambda$ using \cref{corollary:determinant-weight},
the inequalities \eqref{equation:teleman-inequality}
can be summarized by the inequality\vspace*{-1pt}
\begin{equation}
\label{equation:teleman-inequality-for-real}
\mu(\tuple{d}^1)-\mu(\tuple{d}^m)
<
\sum_{1\leq m<n\leq\ell}\left( \mu(\tuple{d}^n)-\mu(\tuple{d}^m) \right)\langle\tuple{d}^n,\tuple{d}^m\rangle.
\end{equation}
This inequality is an entirely algorithmic condition that is implemented in \cite{quivertools},
using the recursion from \cite[Cor.\,3.5]{MR1974891} to enumerate the Harder--Narasimhan strata.

\begin{proof}
We can rewrite $\eta_\lambda$ as\vspace*{-5pt}
\begin{align*}
\eta_\lambda
&= \sum_{1\leq m < n\leq\ell} (k_n - k_m)\langle \tuple{d}^m,\tuple{d}^n \rangle= \sum_{1\leq m < n\leq\ell} \sum_{r = m}^{n-1} (k_{r+1} - k_r) \langle \tuple{d}^m,\tuple{d}^n \rangle \\
&= \sum_{r = 1}^{\ell-1} (k_{r+1} - k_r)\sum_{\substack{m \leq r \\ n > r}} \langle \tuple{d}^m, \tuple{d}^n \rangle= \sum_{r = 1}^{\ell-1} (k_r - k_{r+1})\underbrace{\sum_{\substack{m \leq r \\ n > r}} \bigl( - \langle \tuple{d}^m, \tuple{d}^n \rangle\bigr)}_{=: N_r}.
\end{align*}
To prove the inequality~\eqref{equation:teleman-inequality},
it will be enough to show that~$N_r \geq 1$ for all~$r$,
and that there exists at least one $r$ for which $N_r \geq 2$.
To apply the condition from \cref{definition:strong-ample-stability}, we~notice that each term $N_r$ satisfies the equality
\begin{equation*}
-N_r = \Bigl\langle \sum_{m \leq r} \tuple{d}^m, \sum_{n > r} \tuple{d}^n \Bigr\rangle
= \langle \tuple{e}, \tuple{d} - \tuple{e} \rangle.
\end{equation*}
Since $\mu(\tuple{e}) > \mu(\tuple{d} - \tuple{e})$, we~can conclude that $N_r \geq 2$.
This completes the proof.
\end{proof}

We hope that the strong ample stability in \cref{proposition:inequality}
is not needed for the statement of \cref{theorem:cohomology-vanishing},
and that ample stability in fact suffices
(cf.~\cref{example:strict-ample-stability-rigid}),
or in fact that it holds as soon as there exists a universal representation.
However, the following shows that the use of Teleman quantization,
and thus \cref{proposition:inequality},
does require at least ample stability to hold.

\begin{lemma}
\label{lemma:teleman-implies-ample-stability}
Let~$Q$ be a quiver,
$\tuple{d}$ a dimension vector,
and~$\theta$ a stability parameter,
such that~$\tuple{d}$ is not amply~$\theta$-stable.
Let~$\lambda=\lambda_{\tuple{d}^*}$ be the $1$-parameter subgroup
corresponding to one of the Harder--Narasimhan types~$\tuple{d}^*$
whose associated stratum is of codimension~$1$.
Then \eqref{equation:teleman-inequality-for-real} does not hold.
\end{lemma}

\begin{proof}
By \cite[Prop.\,3.4]{MR1974891},
the codimension of the Harder--Narasimhan stratum can be computed as
\begin{equation}
\label{equation:codimension-one-sum}
1=\sum_{1\leq m<n\leq\ell}-\langle\tuple{d}^m,\tuple{d}^n\rangle
\end{equation}
and because~$m<n$ in this sum, we~have
\begin{equation*}
\langle\tuple{d}^m,\tuple{d}^n\rangle
=\operatorname{hom}(\tuple{d}^m,\tuple{d}^n)-\operatorname{ext}(\tuple{d}^m,\tuple{d}^n)
=-\operatorname{ext}(\tuple{d}^m,\tuple{d}^n).
\end{equation*}
Thus in the sum~\eqref{equation:codimension-one-sum}
all but one term vanishes,
with one term having value exactly~1.
But then
\begin{equation*}
k_\ell-k_1\geq\sum_{1\leq m<n\leq\ell}(k_m-k_n)(-\langle\tuple{d}^m,\tuple{d}^n\rangle)=\eta_\lambda,
\end{equation*}
where the last equality is \cref{corollary:determinant-weight},
thus~\eqref{equation:teleman-inequality-for-real} does not hold.
\end{proof}

We now come to the proof of the cohomology vanishing.

\begin{proof}[Proof of \cref{theorem:cohomology-vanishing}]
In order to apply Teleman quantization
as stated in \cref{theorem:teleman-quantization}
we need to check that,
for every $1$-parameter subgroup~$\lambda$
arising in the Hesselink stratification,
the limit set~$\limitset$ is connected,
and the weights of the equivariant bundle are bounded above by~$\eta_\lambda$ as defined in~\eqref{equation:eta-lambda}.
The connectedness follows from the definition of~$\limitset[{\tuple{d}^*}]$
in \cref{definition:harder-narasimhan-strata}
and the correspondence in \cref{theorem:HN=hesselink}.
That the inequality holds whenever~$\tuple{d}$ is strongly amply~$\theta$-stable
is checked in \cref{proposition:inequality}.

Thus we can apply Teleman quantization,
and we obtain isomorphisms
\begin{equation*}
\HH^k(\Rep(Q,\tuple{d}),U_i^\vee\otimes U_j)^{\group}
\cong
\HH^k(\Rep^{\theta\semistable}(Q,\tuple{d}),U_i^\vee\otimes U_j)^{\group}
\cong
\HH^k(\modulispace[\theta\stable]{Q,\tuple{d}},\mathcal{U}_i^\vee\otimes\mathcal{U}_j)
\end{equation*}
for all~$k\geq 0$, as~stability and semistability agree.
But because the representation variety~$\Rep(Q,\tuple{d})$ is \emph{affine},
the cohomology~$\HH^k(\Rep(Q,\tuple{d}),U_i^\vee\otimes U_j)$ vanishes for~$k\geq 1$,
even before taking invariants,
which proves~\eqref{equation:cohomology-vanishing}.
\end{proof}

As already explained in the introduction,
\eqref{equation:Ext-as-cohomology} shows that \cref{theorem:cohomology-vanishing}
is equivalent to \cref{corollary:schofield}
and thus proves Schofield's conjecture
under the additional assumption of strong ample stability,
without needing that~$Q$ is acyclic.

\subsubsection*{Windows and weights for a $6$-dimensional Kronecker moduli space}
To illustrate the theory we will give the details of an interesting and relevant example.
We will consider the smallest Kronecker moduli space which is not a projective space or a Grassmannian,
whose geometry was studied explicitly in \cite{MR4770368}.
\begin{example}
\label{example:kronecker-moduli-space}
Let~$Q$ be the $3$-Kronecker quiver,
and consider the dimension vector~$\tuple{d}=(2,3)$.
The stability parameter~$\theta=(3,-2)$ is necessarily
(a rescaling of) the canonical stability parameter~$\theta_\can=(9,-6)$.
\end{example}
We have~$\Rep(Q,\tuple{d})\cong\Mat_{3\times 2}(\field)^3\cong\mathbb{A}^{18}$,
with a group action of~$\GL_2\times\GL_3$.
There are~8 Harder--Nara\-sim\-han strata,
the (semi)stable one and~7 unstable strata. In~\cref{table:harder-narasimhan-strata} we have enumerated the properties of the unstable strata.
Everything can be determined from the Euler pairings~$(\langle\tuple{d}^m,\tuple{d}^n\rangle)_{m,n=1,\ldots,\ell}$
and the tuple~$(\mu(\tuple{d}^m))_{m=1,\ldots,\ell}$.
The Harder--Narasimhan type~$((2,0),(0,3))$ is the origin in~$\mathbb{A}^{18}$,
corresponding to the direct sum~$\mathrm{S}_1^{\oplus2}\oplus\mathrm{S}_2^{\oplus3}$ of simple representations.

In \cref{table:harder-narasimhan-strata} we only list the difference~$k_1-k_\ell$, as~explained in the discussion surrounding \eqref{equation:teleman-inequality-for-real}, as~this is sufficient to check the conditions to apply Teleman quantization.

\begin{table}
\centering
\caption{Harder--Narasimhan strata and their properties for 3\protect\nobreakdash-Kro\-necker quiver and $\tuple{d}=(2,3)$.}
\label{table:harder-narasimhan-strata}
\begin{tabular}{lcccccc}
\toprule
\multicolumn{1}{c}{$\tuple{d}^*$} & $\operatorname{codim}$ & $(\mu(\tuple{d}^m))_{m=1,\ldots,\ell}$ & $C$ & $(k_m)_{m=1,\ldots,\ell}$ & $k_1-k_\ell$ & $\eta_\lambda$ \\
\midrule
$((1,1),(1,2))$ & 3 & $(1/2,-1/3)$ & 6 & $(3,-2)$ & 5 & 15 \\
$((2,2),(0,1))$ & 4 & $(1/2,-2)$ & 2 & $(1,-4)$ & 5 & 20 \\
$((2,1),(0,2))$ & 10 & $(4/3,-2)$ & 3 & $(4,-6)$ & 10 & 100 \\
$((1,0),(1,3))$ & 8 & $(3,-3/4)$ & 4 & $(12,-3)$ & 15 & 120 \\
$((1,0),(1,2),(0,1))$ & 9 & $(3,-1/3,-2)$ & 3 & $(9,-1,-6)$ & 15 & 100 \\
$((1,0),(1,1),(0,2))$ & 12 & $(3,1/2,-2)$ & 2 & $(6,1,-4)$ & 10 & 90 \\
$((2,0),(0,3))$ & 18 & $(3,-2)$ & 1 & $(3,-2)$ & 5 & 90 \\
\bottomrule
\end{tabular}
\end{table}

\subsubsection*{Schofield's conjecture in the twisted case}
The statement of Schofield's conjecture uses the universal bundle,
which is thus required to exist. By~\cite[Th.\,4.4]{MR3683503} there are obstructions to its existence
if~$\gcd(\tuple{d})\geq 2$.
However, it is reasonable to expect a \emph{twisted} universal bundle~$\mathcal{U}$,
defined (only) on the stable locus,
which can be seen as an alternative natural generator of the Brauer group in \cite[Conj.\,4.3]{MR3683503}.
In~$\mathcal{U}^\vee\otimes\mathcal{U}$ the twists cancel each other,
and this is in fact an Azumaya algebra whose class in the Brauer group is the same class.

The following example illustrates how \cref{theorem:cohomology-vanishing}
suggests a twisted version of Schofield's conjecture.

\begin{example}
\label{example:twisted-schofield}
Let~$Q$ be the $3$-Kronecker quiver,
and consider the dimension vector~$\tuple{d}=(2,2)$.
The stability parameter~$\theta=(2,-2)$ is necessarily
(a rescaling of) the canonical stability parameter~$\theta_\can=(6,-6)$.
Then by \cite[p.\,1779]{MR1972892}
and \cite[\S7]{MR3683503} we have that~$\moduli^{\theta\semistable}(Q,\tuple{d})\cong\mathbb{P}^5$
and~$\moduli^{\theta\stable}(Q,\tuple{d})$ is the complement of a determinantal cubic hypersurface.

By \cite[Prop.\,7.1]{MR3683503} we have that~$\Br(\moduli^{\theta\stable}(Q,\tuple{d}))\cong\mathbb{Z}/2\mathbb{Z}$,
and there cannot exist a choice of linearization~$\tuple{a}$
for which the universal representation~$U(\tuple{a})$ descends.
However, it should be a \emph{twisted} universal bundle~$\mathcal{U}$ on~$\moduli^{\theta\stable}(Q,\tuple{d})$,
which corresponds to the unique non-trivial Brauer class.

The twisted analogue of \cref{theorem:cohomology-vanishing}
then implies that~$\mathcal{U}$ is a twisted (partial) tilting bundle,
whose endomorphism algebra (which is an untwisted object)
is an Azumaya algebra whose class in the Brauer group coincides with the twist of~$\mathcal{U}$.
Schofield's conjecture thus holds for the \emph{twisted} moduli space of stable representations.
Observe that~$\moduli^{\theta\stable}(Q,\tuple{d})$ is affine, so~any untwisted vector bundle is already a tilting bundle,
and any twisted vector bundle is a twisted tilting bundle,
without the need for the adjective partial.
\end{example}

\subsection{Rigidity}
\label{subsection:rigidity}
From now on we will assume that~$Q$ is acyclic. In~\cref{section:moduli} we recalled the construction of the universal bundle~$\mathcal{U}=\bigoplus_{i\in Q_0}\mathcal{U}_i$
on~$\modulispace[\theta\stable]{Q,\tuple{d}}$.
\hbox{Using} the summands~$\mathcal{U}_i$ we obtain an exact sequence of vector bundles, by~combining \hbox{\cite[Prop.\,3.3 \& Prop.\,3.7]{MR4770368}}:
\begin{equation}
\label{equation:4-term-exact-sequence}
0
\ra\mathcal{O}_{\modulispace[\theta\stable]{Q,\tuple{d}}}
\to\bigoplus_{i \in Q_0} \mathcal{U}^\vee_i \otimes \mathcal{U}_i
\to\bigoplus_{a \in Q_1} \mathcal{U}^\vee_{\source(a)} \otimes \mathcal{U}_{\target(a)}
\to\tangent_{\modulispace[\theta\stable]{Q,\tuple{d}}}
\ra0.
\end{equation}
See also~\cite[\S 4.1]{MR4352662} for a more direct (but less detailed) construction.

To compute the higher cohomology of the tangent bundle for the proof of \cref{corollary:rigidity}, we~can split the sequence in two short exact sequences.
The cohomology of the middle terms in~\eqref{equation:4-term-exact-sequence}
is the subject of \cref{theorem:cohomology-vanishing}.
For the first term we have the following.

\begin{proposition}
Let~$Q$, $\tuple{d}$ and~$\theta$ be as in \cref{corollary:rigidity},
where it is possible to omit the condition that~$\tuple{d}$ is strongly amply~$\theta$-stable.
Then
\begin{equation}
\label{equation:higher-cohomology-structure-sheaf-vanishing}
\HH^k(\modulispace[\theta\stable]{Q,\tuple{d}},\mathcal{O}_{\modulispace[\theta\stable]{Q,\tuple{d}}}) = 0
\end{equation}
for all $k \geq 1$.
\end{proposition}

\begin{proof}
Because~$\tuple{d}$ is chosen to be $\theta$-coprime,
$\tuple{d}$ is a Schur root for the quiver~$Q$.
As~$\gcd(\tuple{d}) = 1$
we have that~$\modulispace[\theta\stable]{Q,\tuple{d}}$ is a rational variety~\cite[Th.\,6.4]{MR1914089}. It~is also smooth and projective, as~discussed in \cref{section:moduli}.
We thus obtain the vanishing in~\eqref{equation:higher-cohomology-structure-sheaf-vanishing}
by the birational invariance of these cohomology groups.
\end{proof}

We thus arrive at the following (short) proof of the rigidity of quiver moduli.

\begin{proof}[Proof of \cref{corollary:rigidity}]
The higher cohomology of the first term in~\eqref{equation:4-term-exact-sequence} vanishes. By~\cref{theorem:cohomology-vanishing} the higher cohomology of the second and third term in~\eqref{equation:4-term-exact-sequence} vanishes,
if we in addition assume \cref{definition:strong-ample-stability}.
This allows us to conclude.
\end{proof}

\begin{example}
\label{example:strict-ample-stability-rigid}
Consider the setup of \cref{example:strict-ample-stability},
which gave an example of a quiver, dimension vector and stability parameter
which was amply stable, but not strongly amply stable.
The inequality in~\eqref{equation:teleman-inequality-for-real} can be shown to hold directly in this case,
thus the resulting moduli space is rigid, even in the absence of strong ample stability.
\end{example}

We now give an example of a moduli space that is rigid
even though the inequality~\eqref{equation:teleman-inequality-for-real} is not satisfied,
and in fact, the defining data is not even amply stable.
\begin{example}
\label{example:no-teleman-still-rigid}
Let $Q$ be the $3$-vertex quiver
\begin{equation}
%\label{equation:3-vertex-quiver-2}
Q\colon
\begin{tikzpicture}[baseline = -20pt, node distance = 1.5cm]
\node (1) {};
\node (2) [right of = 1] {};
\node (3) [below of = 2] {};
\draw (1) circle (2pt) node[above] {1};
\draw (2) circle (2pt) node[above] {2};
\draw (3) circle (2pt) node[below] {3};
\draw[->] (1) edge (2);
\draw[->, bend left = 25] (1) edge (3);
\draw[->, bend left = 15] (1) edge (3);
\draw[->, bend left = 5] (1) edge (3);
\draw[->, bend right = 5] (1) edge (3);
\draw[->, bend right = 15] (1) edge (3);
\draw[->, bend right = 30] (1) edge (3);
\draw[->] (2) edge (3);
\end{tikzpicture}
\end{equation}
Let $\tuple{d} = (1,6,6)$ and let $\theta = \theta_\can = (42,5,-12)$ be the canonical stability parameter.
We have again that~$\tuple{d}$ is~$\theta$\dash coprime, and that the $\theta$-stable locus is non-empty.

The Harder--Narasimhan stratification of the representation variety contains 85~strata:
84~unstable strata, plus the dense stratum of stables.
This can easily be verified using \cite{quivertools}.
The stratum of HN-type $((0,1,0), (1,5,6))$ has codimension~$1$, so~$\tuple{d}$ is not amply~$\theta$-stable.
The condition in~\eqref{equation:teleman-inequality-for-real} does not hold, as~the same stratum of HN-type $((0,1,0), (1,5,6))$ gives~$k_1 = 5$, $k_2 = -\sfrac{5}{12}$,
and~$\langle (0,1,0),(1,5,6) \rangle = -1$.
However, the resulting moduli space is $\mathbb{P}^6$ by \cref{lemma:P6},
for which rigidity easily follows from the Euler sequence.
\end{example}

\begin{lemma}
\label{lemma:P6}
For~$Q,\tuple{d},\theta$ as in \cref{example:no-teleman-still-rigid}
there exists an isomorphism
\begin{equation*}
\modulispace[\theta\stable]{Q,\tuple{d}}\cong\mathbb{P}^6.
\end{equation*}
\end{lemma}

\begin{proof}
The data of a representation $M$ consists of
a vector~$v$ in~$M_2\cong\field^6$,
six vectors~$w_1,\dots,w_6$ in~$M_3\cong\field^6$,
and an endomorphism of~$\field^6$ that we denote by~\hbox{$A \colon \field^6 \!\to\! \field^6$}.
Let us consider the conditions on this data for $M$ to be $\theta$-stable.
If $A$ has a nontrivial kernel, then $M$ admits a subrepresentation of dimension vector $(0,\dim_\field\ker A,0)$ or $(1,\dim_\field\ker A,0)$,
and both have positive slope.
Thus $A$ must be injective for~$M$ to be~$\theta$-stable.
If $Av,w_1,\dots,w_6$ together do not span $M_3$,
then there is a subrepresentation of dimension vector $(1,6,5)$ which has positive slope.
Therefore, the $6 \times 7$-matrix
\begin{equation*}
F_{M} \colonequals\left( Av~|~w_1~|~\dots~|~w_6 \right)
\end{equation*}
must have rank $6$.
Using this description we can see that
the moduli space can be identified to $\mathbb{P}^6$.
To do so, we~send the data of a semistable representation $M$ to
the kernel of the linear map $F_M \colon \field^7 \to \field^6$,
which is a line in~$\field^7$.
\end{proof}
\ifcode
Corresponding Sage code below. Version of quiver.tools: commit a07f4f0

\begin{verbatim}
sage: from quiver import *
sage: Q = ThreeVertexQuiver(1,6,1)
sage: d = vector([1,6,6])
sage: theta = Q.canonical_stability_parameter(d)
sage: theta
(42, 5, -12)
sage: is_coprime_for_stability_parameter(d,theta)
True
sage: Q.has_semistable_representation(d,theta)
True
sage: len(Q.all_harder_narasimhan_types(d,theta))
85
sage: hntype = [vector([0,1,0]),vector([1,5,6])]
sage: hntype in Q.all_harder_narasimhan_types(d,theta)
True
sage: Q.codimension_of_harder_narasimhan_stratum(hntype)
1
sage: k1 = slope(hntype[0],theta)
sage: k2 = slope(hntype[1],theta)
sage: print(k1,k2)
5 -5/12
sage: Q.euler_form(hntype[0],hntype[1])
-1
\end{verbatim}
\fi

\subsection{Height-zero moduli spaces}
\label{subsection:height-zero}
Let us recall some aspects of the moduli theory of semistable sheaves on~$\mathbb{P}^2$.
We will denote
the moduli space of Gieseker-semistable sheaves with given rank and first (\resp second) Chern class
by~$\moduli_{\mathbb{P}^2}(r,\mathrm{c}_1,\mathrm{c}_2)$.
In~\cite{MR0916199} a function~$\delta\colon\mathbb{Q}\to\mathbb{Q}$ is constructed,
for which it is shown that~$\moduli_{\mathbb{P}^2}(r,\mathrm{c}_1,\mathrm{c}_2)$ has strictly positive dimension
if and only if
\begin{equation*}
\delta\left( \sfrac{\mathrm{c}_1}{r} \right)
\leq
\frac{1}{r}\Bigl( \mathrm{c}_2 - \Bigl( 1-\frac{1}{r} \Bigr)\frac{\mathrm{c}_1^2}{2} \Bigr).
\end{equation*}
When it is an \emph{equality},
the moduli space is said to be of \emph{height zero}.

It is also shown in op.~cit.~that for moduli spaces of strictly positive dimension
there exists an associated exceptional vector bundle.
Writing~$\mu\colonequals\sfrac{\mathrm{c}_1}{r}$ for the slope of the sheaves parametrized by~$\moduli_{\mathbb{P}^2}(r,\mathrm{c}_1,\mathrm{c}_2)$, we~will denote this associated exceptional vector bundle by~$E_\mu$.

As explained in~\cite{MR0916199},
moduli spaces of height zero have special properties
amongst all moduli spaces of semistable sheaves on~$\mathbb{P}^2$.
The one which is relevant to us is the following identification~\cite[Th.\,2]{MR0916199}.

\begin{theorem}[Drezet]
\label{theorem:identification}
Let~$\moduli_{\mathbb{P}^2}(r,\mathrm{c}_1,\mathrm{c}_2)$ be a moduli space of height zero.
Then there exists a natural isomorphism
\begin{equation*}
\moduli_{\mathbb{P}^2}(r,\mathrm{c}_1,\mathrm{c}_2)
\cong
\modulispace[\theta_\can\semistable]{\mathrm{K}_{3r_\mu},(m,n)}
\end{equation*}
where~$r_\mu$ is the rank of the of the associated exceptional vector bundle~$E_\mu$,
$\mathrm{K}_{3r_\mu}$ is the Kronecker quiver with~$3r_\mu$ arrows,
and the dimension vector~$(m,n)$ is determined as in~\cite[\S IV.2]{MR0916199}.
\end{theorem}

In what follows, we~only consider invariants~$(r,\mathrm{c}_1,\mathrm{c}_2)$
such that every semistable sheaf is stable,
and thus~$(m,n)$ will be coprime.

For Kronecker quivers,~\cite[Prop.\,6.2]{MR3683503}
shows that strong ample stability holds, as~also recalled in \cref{lemma:strong-ample-stability-examples}.
We thus have the following observation which we record for later use.

\begin{lemma}
\label{lemma:rigidity-kronecker-quivers}
Let $Q$ be a Kronecker quiver,
and let~$\tuple{d}$ be a dimension vector
which is~$\theta_\can$-coprime (and thus indivisible).
Then~$\modulispace[\theta\semistable]{Q,\tuple{d}}$ is rigid.
\end{lemma}

To prove \cref{corollary:height-zero}, we~wish to obtain a contradiction on the assumption that~\eqref{equation:height-zero-functor} is fully faithful.
The combination of~\cite[Prop.\,29, Rem.\,30]{MR3950704} states the following
for fully faithful functors towards moduli spaces of sheaves on surfaces.
\begin{theorem}[Belmans--Fu--Raedschelders]
\label{theorem:fully-faithful-hochschild-cohomology}
Let~$S$ be a smooth projective surface
such that~$\mathcal{O}_S$ is exceptional,
i.e.,~$\HH^1(S,\mathcal{O}_S)=\HH^2(S,\mathcal{O}_S)=0$.
Let~$\moduli_S(r,\mathrm{c}_1,\mathrm{c}_2)$ be a smooth projective moduli space
of stable sheaves on~$S$ with prescribed invariants.
Let~$\mathcal{E}$ be the universal sheaf.
Assume that the Fourier--Mukai functor
\begin{equation*}
\Phi_{\mathcal{E}}\colon\derived^\bounded(S)\to\derived^\bounded(\moduli_S(r,\mathrm{c}_1,\mathrm{c}_2))
\end{equation*}
is fully faithful.
Then
\begin{equation*}
\HH^i(\moduli_S(r,\mathrm{c}_1,\mathrm{c}_2),\tangent_{\moduli_S(r,\mathrm{c}_1,\mathrm{c}_2)})
\cong
\hochschild^{i+1}(S).
\end{equation*}
\end{theorem}

\begin{proof}[Proof of \cref{corollary:height-zero}]
For~$\mathbb{P}^2$,
the Hochschild--Kostant--Rosenberg decomposition
gives us the isomorphism~$\hochschild^2(\mathbb{P}^2)\cong\HH^0(\mathbb{P}^2,\omega_{\mathbb{P}^2}^\vee)\cong\field^{10}$. By~the identification from \cref{theorem:identification}
and the assumption that the moduli spaces are smooth
we have that
\begin{equation}
\label{equation:first-cohomology-tangent-bundle-identification}
\HH^1(\moduli_{\mathbb{P}^2}(r,\mathrm{c}_1,\mathrm{c}_2),\tangent_{\moduli_{\mathbb{P}^2}(r,\mathrm{c}_1,\mathrm{c}_2)})
\cong
\HH^1(\modulispace[\theta_\can]{\mathrm{K}_{3r_\mu},(m,n)},\tangent_{\modulispace[\theta_\can]{\mathrm{K}_{3r_\mu},(m,n)}}).
\end{equation}
But by \cref{corollary:rigidity},
which we can apply because of \cref{lemma:rigidity-kronecker-quivers},
the cohomology in~\eqref{equation:first-cohomology-tangent-bundle-identification} vanishes,
which contradicts \cref{theorem:fully-faithful-hochschild-cohomology}.
\end{proof}

The following remark gives an alternative method to prove
the functor cannot be fully faithful,
using techniques introduced already by Drezet.
\begin{remark}
\label{remark:associated-exceptional-bundle}
Consider~$\moduli_{\mathbb{P}^2}(r,\mathrm{c}_1,\mathrm{c}_2)$ of height zero.
As explained in~\cite[\S5.1]{MR0916199},
there exists a unique associated exceptional vector bundle~$E$
for which
\begin{equation}
\label{equation:inequalities}
-3<\frac{\mathrm{c}_1H}{r}+\frac{\mathrm{c}_1(E)H}{\rk E}\leq 0.
\end{equation}
As mentioned in loc.~cit.,
this implies that
\begin{equation*}
\chi(\mathbb{P}^2,E^\vee\otimes V)=\operatorname{ht}(\moduli_{\mathbb{P}^2}(r,\mathrm{c}_1,\mathrm{c}_2))=0
\end{equation*}
for all~$[V]\in\moduli_{\mathbb{P}^2}(r,\mathrm{c}_1,\mathrm{c}_2)$.
But then we can prove that~$\Phi_{\mathcal{U}}(E^\vee)=0$,
contradicting fully-faithfulness.
Namely, fiberwise we see for the transform of~$E^\vee$ that
\begin{equation*}
\HH^\csbullet(\mathbb{P}^2,(p^*E^\vee\otimes\mathcal{U})_{\mathbb{P}^2\times[V]})
\cong
\HH^\csbullet(\mathbb{P}^2,E^\vee\otimes V)
=
0,
\end{equation*}
as both~$\Hom(E,V)$ and~$\Ext^2(E,V)=\Hom(V,E\otimes\mathcal{O}_{\mathbb{P}^2}(3))^\vee$ vanish, by~the inequalities~\eqref{equation:inequalities}.
\end{remark}

\appendix

\section*{Appendix. Proof of \texorpdfstring{\cref{proposition:strong-ample-stability-fundamental-domain}}{Proposition \ref{proposition:strong-ample-stability-fundamental-domain}}}
\refstepcounter{section}
\label{section:proof-no-thin-bridges}
We now come to the proof of \cref{proposition:strong-ample-stability-fundamental-domain},
which we relegated to this appendix because it is rather lengthy
and the methods are not related to the rest of the paper.

Recall that a quiver~$Q$ and a dimension vector~$\tuple{d}$
have a \emph{thin bridge}
if there exists a decomposition~$Q_0=I'\sqcup I''$
with a unique arrow~$a\colon i'\to i''$ connecting~$I'$ and~$I''$
(so~$i'\in I'$ and~$i''\in I''$)
so that~$d_{i'}=d_{i''}=1$.

\begin{proof}[Proof of \cref{proposition:strong-ample-stability-fundamental-domain}]
Suppose that~$\tuple{d}$ is not strongly amply~$\theta_\can$-stable, so~that
there exists a subdimension vector~$\tuple{e}\leq\tuple{d}$
such that~$\theta_\can(\tuple{e})\geq 1$
and~$\langle\tuple{e},\tuple{d}-\tuple{e}\rangle\geq-1$.
This yields
\begin{equation*}
\{\tuple{e},\tuple{d}-\tuple{e}\}
=\{\tuple{e},\tuple{d}\}
=-\{\tuple{d},\tuple{e}\}
=-\theta_\can(\tuple{e})\leq-1,
\end{equation*}
which implies that
%\begin{equation*}
$\langle\tuple{d}-\tuple{e},\tuple{e}\rangle
=\langle\tuple{e},\tuple{d}-\tuple{e}\rangle-\{\tuple{e},\tuple{d}-\tuple{e}\}
\geq0$,
%\end{equation*}
and thus
\begin{equation*}
(\tuple{d}-\tuple{e},\tuple{e})\geq-1.
\end{equation*}
As in the proof of \cite[Lem.\,2.15]{MR0557581},
and more precisely Equation~(2.10) of op.~cit., we~can write
\begin{equation}
\label{equation:kac-inequality}
-1
\leq({\tuple{d}}-{\tuple{e}},{\tuple{e}})
=\underbrace{\sum_{i\in Q_0}\frac{e_i(d_i-e_i)}{d_i}\,({\tuple{d}},{\tuple{i}})}_{S_1}
+\underbrace{\frac{1}{2}\sum_{i\neq j\in Q_0}d_id_j\Bigl( \frac{e_i}{d_i}-\frac{e_j}{d_j} \Bigr)^2({\tuple{i}},{\tuple{j}})}_{S_2}.
\end{equation}
Since~$\tuple{d}$ is assumed in the interior of the fundamental domain,
the summand~$S_1$ is non-positive,
whereas~$S_2$ is non-positive because~$(\tuple{i},\tuple{j})$ is always non-positive.

First, assume that~$S_2=0$.
This implies that the vectors~$(e_i,d_i)$ and~$(e_j,d_j)$ are proportional
whenever~$i$ and~$j$ are connected in~$Q$.
Because~$\tuple{d}$ is~$\theta_\can$-coprime, the quiver is connected.
But then the proportionality implies that~$\tuple{e}$ is proportional to~$\tuple{d}$,
which contradicts the~$\theta_\can$-coprimality.
Hence~$S_2<0$.

Next, we~will consider~$S_1$.
We want to show that~$e_i \in \{0,d_i\}$ for all~$i\in Q_0$.
If there exists some vertex~$k\in Q_0$
for which~$e_k\not\in\{0,d_k\}$,
then~$d_k\geq 2$,
and we have the bound
\begin{equation}
\label{equation:two-inequalities}
\frac{e_k(d_k-e_k)}{d_k}\geq 1-\frac{1}{d_k}\geq\frac{1}{2}
\end{equation}
on the coefficient of~$(\tuple{d},\tuple{k})$ in the~$k$th term of~$S_1$.

If there are at least two vertices for which~$e_k\not\in \{0,d_k\}$,
then
\begin{equation*}
S_1\leq-\frac{1}{2}-\frac{1}{2}=-1,
\end{equation*}
because we have taken~$\tuple{d}$ in the interior of the fundamental domain.
However, this contradicts \eqref{equation:kac-inequality},
as~$S_2<0$.
Therefore, there is at most one vertex $k$ of the quiver
for which~$e_k \notin \{0,d_k\}$.

Assume now that there is precisely one such vertex $k$.
We have~$e_k\in\{1,\ldots,d_k-1\}$, which using \eqref{equation:two-inequalities}
implies $S_1 < -(1 - \sfrac{1}{d_k})$.

Let us see what this does to~$S_2$.
Since $Q$ is connected, there exists at least one vertex~$j$ of~$Q$
which is connected to $k$ by at least one arrow,
which means that there are at least two summands in $S_2$ that are nonzero.
As $j \neq k$, $e_j \in \{0, d_j\}$ by hypothesis, so~that the respective summands simplify to~$d_j\frac{(d_k-e_k)^2}{d_k}(\tuple{k},\tuple{j})$
or~$d_j\frac{e_k^2}{d_k}(\tuple{k},\tuple{j})$. In~either case, their value is bounded above by~$-\sfrac{1}{d_k}$, so~that~$S_2\leq-\sfrac{1}{d_k}$.
If this inequality is strict,
then together with the first inequality in \eqref{equation:two-inequalities}
it contradicts \eqref{equation:kac-inequality},
thus~$S_1=0$.

We will show now that the inequality above is in fact strict, i.e., that $S_2 < \sfrac{1}{d_k}$.
This amounts to proving that there is at least another summand in $S_2$ that is nonzero.

Let us split the quiver $Q$ in three disjoint sets of vertices,
i.e., we~decompose~$Q_0$ into
the disjoint union~$\{k\} \sqcup I' \sqcup I''$,
where~$I' \colonequals \{j \mid e_j = d_j\}$ and~$I'' \colonequals \{j \mid e_j = 0\}$.
We see that any arrow from $I'$ to $I''$ or vice versa yields a nonzero summand of~$S_2$, so~to conclude we must show that neither $I'$ nor $I''$ are empty.
As $Q$ is connected, this will imply that at least one such arrow exists.

Let us assume by contradiction that no such arrow exists, i.e.,
that in the previous decomposition either of $I'$ or $I''$ is empty.
This means that $S_2 = \sfrac{1}{d_k}$, and the inequality \eqref{equation:kac-inequality}
yields $-1 \leq S_1 + S_2 \leq -1$, so~that~$(\tuple{d} - \tuple{e}, \tuple{e}) = -1$.
This gives the two following identities:
\begin{equation}\label{equation:contradiction-1}
\langle \tuple{e}, \tuple{d} - \tuple{e} \rangle = -1
\end{equation}
and
\begin{equation}\label{equation:contradiction-2}
\langle \tuple{d} - \tuple{e}, \tuple{e} \rangle = 0.
\end{equation}

These two identities contradict the fact that $\tuple{d}$ belongs to the fundamental domain.
We will show this for the case $I' = \emptyset$, the other case being essentially the same.

Using the conditions on~$e_i$ for~$i\neq k$,
the equations \eqref{equation:contradiction-1} and \eqref{equation:contradiction-2} are written out explicitly as\vspace*{-3pt}
\begin{equation*}
\begin{aligned}
\langle \tuple{e}, \tuple{d} - \tuple{e} \rangle = -1
&= e_k(d_k - e_k) - \sum_{a \in Q_1}e_{\source(a)}(d_{\target(a)} - e_{\target(a)}) \\
&= e_k(d_k - e_k) - \sum_{\substack{a\in Q_1 \\ \source(a)=k}} e_{k} d_{\target(a)}
\end{aligned}
\end{equation*}
and\vspace*{-3pt}
\begin{equation*}
\begin{aligned}
\langle \tuple{d} - \tuple{e}, \tuple{e} \rangle = 0
&= e_k(d_k - e_k) - \sum_{a \in Q_1}(d_{\source(a)} - e_{\source(a)})e_{\target(a)} \\
&= e_k(d_k - e_k) - \sum_{\substack{a\in Q_1 \\ \target(a)=k}} d_{\source(a)} e_{k}.
\end{aligned}
\end{equation*}
These yield the inequalities\vspace*{-3pt}
\begin{equation*}
\sum_{\substack{a\in Q_1 \\ \target(a)=k}} d_{\source(a)} = d_k - e_k < d_k, \qquad
\sum_{\substack{a\in Q_1 \\ \source(a)=k}} d_{\target(a)} = d_k - e_k + \frac{1}{e_k} \leq d_k.
\end{equation*}
Summing them we obtain\vspace*{-3pt}
\begin{equation*}
2d_k - \sum_{\substack{a\in Q_1 \\ \source(a)=k}} d_{\target(a)} - \sum_{\substack{a\in Q_1\\\target(a)=k}} d_{\source(a)} > 0.
\end{equation*}
The left-hand side is $(\tuple{d},\tuple{k})$, and
this inequality contradicts the assumption that $(\tuple{d},\tuple{k}) \leq -1$.
We conclude that there cannot be a vertex $k \in Q$ for which $e_k \notin \{0,d_k\}$.
Thus, $S_1 = 0$.

We obtain the equality~$(\tuple{d}-\tuple{e},\tuple{e})=-1$ in \eqref{equation:kac-inequality},
and thus\vspace*{-3pt}
\begin{equation}
\label{equation:equality-1}
\langle\tuple{e},\tuple{d}-\tuple{e}\rangle=-1
\end{equation}
and\vspace*{-3pt}
\begin{equation}
\label{equation:equality-2}
\langle\tuple{d}-\tuple{e},\tuple{e}\rangle=0.
\end{equation}
Denote again by~$I'\subseteq Q_0$ the set of vertices for which~$e_i=d_i$,
and let~$I''$ be its complement.
From \eqref{equation:equality-1} and \eqref{equation:equality-2}
we obtain that there exists a unique edge~$i'\to i''$ from~$I'$ to~$I''$.
Because~$e_i\in\{0,d_i\}$ we have that every factor~$(\sfrac{e_i}{d_i}-\sfrac{e_j}{d_j})^2$
in~$S_2$ is either~0 or~1.
This implies that~$S_2=-1$ has a unique non-zero summand,
equal to~$-1$,
and thus~$d_{i'}=d_{i''}=1$.
We obtain a thin bridge,
and this contradiction finishes the proof.
\end{proof}

\backmatter
\bibliographystyle{jepplain+eid}
\bibliography{belmans-et-al}
\end{document}
